\documentclass[10pt]{amsart}

% Package Being Used:
\usepackage{amsmath}
\usepackage{amssymb}
\usepackage{bm}
\usepackage{graphicx}
\usepackage{psfrag}
\usepackage{color}
\usepackage{xcolor}
\usepackage{float}
\usepackage{breqn}

% colors
\definecolor{amethyst}{HTML}{8b5cf6}
\definecolor{darren}{HTML}{8b5cf6}
\long\def\darren#1{\textcolor{darren}{#1}}

\definecolor{orange}{rgb}{0.988, 0.416, 0.012}
\definecolor{neo}{rgb}{0.988, 0.416, 0.012}
\long\def\neo#1{\textcolor{neo}{#1}}

\definecolor{brown}{rgb}{0.588, 0.337, 0.208}
\definecolor{grant}{rgb}{0.588, 0.337, 0.208}
\long\def\grant#1{\textcolor{grant}{#1}}

\definecolor{red}{rgb}{1,0,0}
\definecolor{felix}{rgb}{1,0,0}
\long\def\felix#1{\textcolor{felix}{#1}}

\usepackage{hyperref}
\definecolor{MITred}{RGB}{163,31,52}
\hypersetup{
  colorlinks=true,
  linkcolor=MITred,
  citecolor=MITred,
  urlcolor=MITred
}
\usepackage{url}
\usepackage{algpseudocode}
\usepackage{fancyhdr}
\usepackage{mathtools}
\usepackage{tikz-cd}
\usepackage{xy}
\input xy
\xyoption{all}
\usepackage{stmaryrd}
\usepackage{calrsfs}

% Paper Format and Geometry:
\voffset=-1.4mm
\oddsidemargin=14pt
\evensidemargin=14pt
\topmargin=26pt
\headheight=12pt
\textheight=576pt
\textwidth=441pt
\parskip=0pt plus 4pt

% Head Labels:
\pagestyle{fancy}
\fancyhf{}
\renewcommand{\headrulewidth}{0pt}
\renewcommand{\footrulewidth}{0pt}
\fancyhead[CE]{\fontsize{9}{11}\selectfont G. BLITZ, F. GOTTI, D. HAN, \& H. LIANG}
\fancyhead[CO]{\fontsize{9}{11}\selectfont MAXIMAL COMMON DIVISORS IN MONOID ALGEBRAS}
\fancyhead[LE,RO]{\thepage}
\setlength{\headheight}{12pt}

% Theorems-like Format and Numbering:

\newtheorem{theorem}{Theorem}[section]
\newtheorem{lemma}[theorem]{Lemma}
\newtheorem{proposition}[theorem]{Proposition}
\newtheorem{corollary}[theorem]{Corollary}

\theoremstyle{definition}
\newtheorem{definition}[theorem]{Definition}
\newtheorem{example}[theorem]{Example}
\newtheorem{problem}[theorem]{Problem}
\newtheorem{question}[theorem]{Question}

\theoremstyle{remark}
\newtheorem{remark}[theorem]{Remark}

\numberwithin{equation}{section}

% Personalized Commands:
\newcommand{\aaa}{\mathbb{A}}
\newcommand{\cc}{\mathbb{C}}
\newcommand{\ff}{\mathbb{F}}
\newcommand{\nn}{\mathbb{N}}
\newcommand{\pp}{\mathbb{P}}
\newcommand{\ppp}{\mathcal{P}}
\newcommand{\qq}{\mathbb{Q}}
\newcommand{\rr}{\mathbb{R}}
\newcommand{\zz}{\mathbb{Z}}
\newcommand{\ch}{\mathrm{char}}
\newcommand{\lcm}{\mathrm{lcm}}
\newcommand{\mcd}{\mathrm{mcd}}
\providecommand\ldb{\llbracket}
\providecommand\rdb{\rrbracket}
\newcommand\pval{\mathsf{v}_p}
\newcommand{\gp}{\mathcal{G}}
\newcommand{\qf}{\text{qf}}
\newcommand{\ord}{\mathrm{ord}}
\newcommand{\ii}{\mathcal{Irr}}
\newcommand{\uu}{\mathcal{U}}
\newcommand{\ppf}{\mathcal{P}_{\text{fin}}}
\newcommand{\ppx}{\mathcal{P}_{\text{fin}, 0}}
\newcommand{\inter}{\mathsf{int}}
\newcommand{\cone}{\mathsf{cone}}
\newcommand{\norm}[1]{\left\lVert#1\right\rVert}

\newcommand{\gl}{\mathfrak{GL}}
\newcommand{\psp}{\mathfrak{PSP}}

\newcommand{\Pfin}{\mathcal{P}_{\mathrm{fin}}}
\newcommand{\Pfino}{\mathcal{P}_{\mathrm{fin},0}}
\newcommand{\Pfinu}{\mathcal{P}_{\text{fin}, \uu}}
\newcommand{\s}{\mathscr{S}}

\newcommand{\fl}[1]{\lfloor#1\rfloor}
\newcommand{\cl}[1]{\lceil #1\rceil}
\newcommand{\ds}{\displaystyle}
\newcommand{\llrr}[1]{\langle \langle #1 \rangle \rangle}
\DeclareMathOperator{\rank}{rank}

\keywords{monoid algebra, divisibility, maximal common divisor, MCD domain, MCD-finite domain, atomic domain, q-GCD domain}

\subjclass[2020]{Primary: 13A05, 13F15, 20M25; Secondary: 13G05, 20M13, 06F05}

%Primary
%13A05: Divisibility and factorizations in commutative rings
%13F15: Commutative rings defined by factorization properties (e.g., atomic, factorial, half-factorial)
%20M25: Semigroup rings, multiplicative semigroups of rings 

%Secondary
%13G05: Integral domains
%20M13: Arithmetic theory of semigroups
%06F05: Ordered semigroups and monoids


\begin{document}

\mbox{}
\title{Maximal common divisors in monoid algebras}

\author{Grant Blitz}
\address{PRIMES-USA\\MIT\\Cambridge, MA 02139}
\email{grantstar16@gmail.com}

\author{Felix Gotti}
\address{Department of Mathematics\\MIT\\Cambridge, MA 02139}
\email{fgotti@mit.edu}

\author{Darren Han}
\address{Department of Mathematics\\MIT\\Cambridge, MA 02139}
\email{drh@mit.edu}

\author{Hengrui Liang}
\address{PRIMES-USA\\MIT\\Cambridge, MA 02139}
\email{spaceblastxy1@gmail.com}

\date{\today}

\begin{abstract}
  We say that a commutative monoid $M$ has the MCD property if every nonempty finite subset of $M$ has a maximal common divisor (MCD), and we say that $M$ has the MCD-finite property if every nonempty finite subset of $M$ has only finitely many MCDs up to associates. It is well known that every monoid that satisfies the ascending chain condition on principal ideals is an MCD monoid, while every finite factorization monoid is an MCD-finite monoid. In this paper, we study the MCD and MCD-finite properties in the setting of monoid algebras. After identifying a class of rank-$1$ torsion-free MCD monoids, we investigate the ascent of the MCD property to monoid algebras over fields, proving that if a pre-Schreier monoid has the MCD property then its monoid algebras over any field also have the MCD property. Then we prove that, unlike for the case of polynomial extensions, the property of being atomic does not ascend to monoid algebras over fields even when restricted to the class of MCD monoids. In the second part of the paper, we first identify a class of rank-$1$ torsion-free MCD-finite monoids. Then we establish the ascent of the MCD-finite property to polynomial extensions. We conclude the paper proving that the q-GCD property (i.e., the condition that every nonempty finite subset has at most one MCD), which is a condition stronger than the MCD-finite property, does not ascend to monoid algebras over fields even when restricted to the class of rank-$1$ torsion-free monoids.
\end{abstract}



\bigskip
\maketitle







\bigskip
%%%%%%%%%%%%%%%%%%%%%%
%%%%%%%%%%%%%%%%%%%%%%
\section{Introduction}
\label{sec:intro}

The ascent of algebraic properties to polynomial extensions is a central problem in commutative ring theory, as polynomial rings are among the most fundamental ring extensions. Indeed, several milestones in the development of modern algebra concern the ascent of ideal-theoretic or factorization properties to polynomial extensions. For instance, in 1801, Carl Gauss~\cite{cG1801} established that unique factorization ascends to polynomial extensions; equivalently, the polynomial ring over a unique factorization domain (UFD) is again a UFD. Another important historical example is Hilbert's Basis Theorem, established by David Hilbert~\cite{dH1890} in 1890, which states that the Noetherian property (i.e., the condition that every ascending chain of ideals stabilizes) ascends to polynomial extensions of commutative rings. Hilbert's proof is a landmark in the emergence of modern abstract algebra because it helped shift the mathematical paradigm away from exhaustive computation and toward the study of structural properties.
\smallskip

Given a commutative ring $R$ and a commutative monoid $M$, the monoid algebra $R[M]$ of $M$ with coefficients in $R$ is a natural generalization of the classical polynomial ring $R[x]$: its elements are polynomial-like expressions whose supports are finite subsets of $M$. In particular, $R[\nn_0]$ is naturally isomorphic to $R[x]$. We say that a property ascends to monoid algebras over a field $F$ if, whenever a monoid $M$ has the property, the multiplicative monoid of nonzero elements of $F[M]$ also has the property. The MCD and MCD-finite conditions are the two central properties considered in this paper: they concern, respectively, the existence and finiteness of maximal common divisors of finite subsets of commutative monoids and integral domains. The primary purpose of this paper is to study these divisibility properties in monoid algebras, with particular emphasis on their ascent to polynomial extensions and to monoid algebras over fields.
\smallskip

Let $M$ be a commutative cancellative monoid. We say that $M$ is a unique factorization monoid (UFM) if it satisfies the statement of the Fundamental Theorem of Arithmetic, while $M$ is called atomic if it satisfies the existential part of that statement. In a UFM, every nonempty finite subset has a greatest common divisor (GCD). Motivated by the existence of GCDs, we can define a class of monoids that contains all UFMs: following Kaplansky~\cite{iK70}, we say that $M$ is a GCD monoid if every nonempty finite subset of $M$ has a GCD. An integral domain is called a GCD domain if its multiplicative monoid of nonzero elements is a GCD monoid. The class of GCD domains contains all B\'ezout domains and UFDs. Given the importance of GCDs, both GCD monoids and GCD domains have been extensively studied (see~\cite{dA00,LP20} and the references therein).
\smallskip

A common divisor $d \in M$ of a nonempty finite subset $S$ in $M$ is called a maximal common divisor (MCD) of $S$ in $M$ provided that the only common divisors of the subset
\[
  S/d := \Big\{\frac{s}d : s \in S \Big\}
\]
are the units of $M$. The monoid $M$ is an MCD monoid (or has the MCD property) if every nonempty finite subset of $M$ has an MCD. An integral domain is an MCD domain (or has the MCD property) if its multiplicative monoid of nonzero elements is an MCD monoid. One can readily check that every monoid or integral domain satisfying the ascending chain condition on principal ideals (ACCP) has the MCD property.
\smallskip

Roitman studied the MCD property in 1993 in connection with atomicity in polynomial extensions. The ascent problem for atomicity was raised by Gilmer~\cite[page~189]{rG84} in the setting of monoid algebras and later posed for polynomial extensions by Anderson, Anderson, and Zafrullah~\cite[Question~2]{AAZ90}. Roitman~\cite{mR93} answered the latter question negatively by constructing an atomic domain whose polynomial ring is not atomic. In the same paper, he proved that the MCD property provides a natural sufficient condition for atomicity to ascend to polynomial extensions. More recently, the MCD property has been considered by Dani et al.~\cite{DGHLS25} in the setting of commutative power monoids, by Liang, Wang, and Zhong~\cite{LWZ24} in the setting of rank-$1$ torsion-free commutative monoids, and by Gotti and Polo~\cite{GP23} in the setting of semidomains.
\smallskip

In Section~\ref{sec:the MCD property}, we study the MCD property. We first identify two classes of rank-$1$ torsion-free MCD monoids, one of which generalizes a class considered in~\cite[Section~5]{LWZ24}. We then investigate the ascent of the MCD property to monoid algebras over fields. We prove that if a pre-Schreier monoid $M$ has the MCD property, then $F[M]$ has the MCD property for every field $F$. Thus, the MCD property ascends to monoid algebras over fields within the class of pre-Schreier monoids (i.e., commutative monoids whose elements are primal). We then consider the ascent of atomicity to monoid algebras over fields. Motivated by a construction used in~\cite{CG19} and refined in~\cite{GGP25}, we prove that, unlike in the case of polynomial extensions, atomicity does not ascend to monoid algebras over fields even within the class of MCD monoids. It remains open whether, for a fixed field $F$, atomicity ascends to monoid algebras with coefficients in $F$ within the class of MCD monoids.
\smallskip

We say that $M$ is an MCD-finite monoid (or has the MCD-finite property) if every nonempty finite subset of $M$ has only finitely many MCDs up to associates. An integral domain is an MCD-finite domain (or has the MCD-finite property) if its multiplicative monoid of nonzero elements is an MCD-finite monoid. The class of MCD-finite monoids contains two important subclasses: pre-Schreier monoids and finite factorization monoids (i.e., atomic monoids in which every element has only finitely many factorizations).
\smallskip

The term MCD-finite was introduced, and the corresponding property first studied, by Eftekhari and Khorsandi~\cite{EK18} in connection with the ascent of the IDF property. They proved that, for an IDF domain, MCD-finiteness characterizes the ascent of the IDF property to polynomial extensions. An integral domain has the IDF property if every nonzero element has only finitely many irreducible divisors up to associates. The question of whether the IDF property ascends to polynomial extensions was posed in~\cite[Question~2]{AAZ90} and answered negatively by Malcolmson and Okoh~\cite{MO09}, who also proved that it ascends within the class of GCD domains. The ascent of the IDF property has also been considered more recently in~\cite{GZ23,GP25}. Du and Gotti~\cite{DG26} recently studied the MCD-finite property in connection with the unrestricted finite factorization property.
\smallskip

In Section~\ref{sec:the MCD-finite property}, we investigate the MCD-finite property. We first revisit the prime reciprocal monoids introduced in Section~\ref{sec:the MCD property} and prove that they have the MCD-finite property, even though they are not finite factorization monoids and do not satisfy the ACCP. We then establish that the MCD-finite property ascends to polynomial extensions. We conclude the section by considering the q-GCD property, which strengthens the MCD-finite property: a monoid is q-GCD if every nonempty finite subset has at most one MCD up to associates. We prove that the q-GCD property does not ascend to monoid algebras over fields. More precisely, for every prime $p$, we construct a rank-$1$ cancellative torsion-free q-GCD monoid $M$ such that $\ff_p[M]$ is not a q-GCD domain.







\bigskip
%%%%%%%%%%%%%%%%%%%%
%%%%%%%%%%%%%%%%%%%%
\section{Background}
\label{sec:background}

\smallskip
%%%%%%%%%%%%%%%%
\subsection{General Notation}

As is customary, $\zz$, $\qq$, $\rr$, $\aaa$, and $\cc$ will denote the sets of integers, rational numbers, real numbers, algebraic complex numbers, and complex numbers, respectively. We let $\nn$ and $\nn_0$ denote the sets of positive and nonnegative integers, respectively. In addition, we let $\pp$ denote the set of primes. For $p \in \pp$ and $n \in \nn$, we let $\ff_{p^n}$ be the finite field of cardinality $p^n$. For $a,b \in \zz$ with $a \le b$, we let $\ldb a,b \rdb$ denote the set of integers between $a$ and $b$, i.e.,
\[
  \ldb a,b \rdb := \{n \in \zz : a \le n \le b\}.
\]
In addition, for $S \subseteq \rr$ and $r \in \rr$, we set $S_{\ge r} := \{s \in S : s \ge r\}$ and $S_{> r} := \{s \in S : s > r\}$.


\medskip
%%%%%%%%%%%%%%%%%%%
\subsection{Commutative Monoids}

A multiplicatively written commutative semigroup $S$ is called \emph{cancellative} if for all $a,b,c \in S$ the equality $ab = ac$ implies that $b=c$. Although a monoid is usually defined to be a semigroup with an identity element, for the purposes of this paper the term \emph{monoid} refers to a cancellative and commutative semigroup with an identity element.
\smallskip

Let $M$ be a multiplicatively written monoid. We let $M^\bullet := M \setminus \{1\}$. The \emph{group of units} of~$M$ is the abelian group consisting of all invertible elements of~$M$, which is denoted by either $M^\times$ or $\uu(M)$. We say that $M$ is \emph{reduced} if the group of units of $M$ is trivial. The quotient monoid $M/M^\times$ is a reduced monoid, which is denoted by $M_{\text{red}}$ and called the \emph{reduced monoid} of $M$. If $M$ is reduced then we can naturally identify each element $a \in M$ with its class $aM^\times$ in $M_{\text{red}}$.
\smallskip

As the monoid $M$ is cancellative, it can be minimally embedded into an abelian group $\gp(M)$, which is called its \emph{Grothendieck group}: $\gp(M)$ is the abelian group (unique up to isomorphism) such that any abelian group containing an isomorphic copy of $M$ also contains an isomorphic copy of $\gp(M)$. The \emph{rank} of $M$ is defined to be the dimension of the $\qq$-vector space $\qq \otimes_\zz \gp(M)$ (clearly the rank of~$M$ equals the rank of $\gp(M)$). It follows from~\cite[Section~24]{lF70} and~\cite[Theorem~2.9]{rG84} that if $M$ is torsion-free (and cancellative) then $M$ has rank $1$ if and only if it is isomorphic to an additive submonoid of $\qq$. Additive submonoids of the nonnegative cone of $\qq$ have been actively investigated during the last decade: see~\cite{CGGP25} for a recent study of their atomic structure and factorization theory, and see~\cite{BGLZ24,GG25,GGP25,GR25} for recent studies of their monoid algebras.

\medskip
%%%%%%%%%%%%%
\subsection{Divisibility}

For $a,b \in M$, we say that $b$ \emph{divides} $a$ and write $b \mid_M a$ if there exists $c \in M$ such that $a = bc$. For $b,c \in M$, if both relations $b \mid_M c$ and $c \mid_M b$ hold then $b$ and $c$ are called \emph{associates}. Let $S$ be a nonempty subset of $M$. An element $d \in M$ is called a \emph{common divisor} of $S$ if $d \mid_M s$ for all $s \in S$. A common divisor $d \in M$ is called a \emph{greatest common divisor} (GCD) of $S$ if every other common divisor of $S$ divides $d$: we let $\gcd_M(S)$ denote the set of all GCDs of $S$, dropping the subscript $M$ when we see no risk of ambiguity (observe that any two GCDs of $S$ are clearly associates).
\smallskip

A common divisor $d \in M$ of $S$ is called a \emph{maximal common divisor} (MCD) of $S$ if the only common divisors of the set $S/d := \{s/d : s \in S\}$ are the units of~$M$. We let $\mcd_M(S)$ denote the set consisting of all the MCDs of $S$ in $M$, dropping the subscript $M$ when we see no risk of ambiguity. If every nonempty finite subset of $M$ has an MCD, we say that $M$ is an \emph{MCD monoid}. Following~\cite{EK18}, we say that $M$ is an \emph{MCD-finite monoid} if every nonempty finite subset of~$M$ has only finitely many MCDs up to associates. When a monoid is MCD-finite, we also say that it has the \emph{MCD-finite property}. If every nonempty finite subset of $M$ has at most one MCD up to associates then we say that $M$ is a \emph{q-GCD monoid} (or a \emph{quasi-GCD monoid}). It follows directly from the corresponding definitions that every q-GCD monoid has the MCD-finite property.
\smallskip

The monoid $M$ is called a \emph{pre-Schreier monoid} if each element $a \in M$ is \emph{primal}, which means that, for all $b,c \in M$ such that $a \mid_M bc$, we can write $a = b' c'$ for some elements $b', c' \in M$ such that $b' \mid_M b$ and $c' \mid_M c$. Clearly, every unit is primal. If a monoid is pre-Schreier then we say that it has the \emph{pre-Schreier property}. A monoid has the \emph{GL property} if for every $a \in M$ the following condition holds: if $a \mid_M bc$ for some $b,c \in M$ then there exists $d \in M \setminus M^\times$ with $d \mid_M a$ such that $d \mid_M b$ or $d \mid_M c$.



\medskip
%%%%%%%%%%%%%%%%%%%%%%%%%%%%%%%%%%%%%%%%
\subsection{Atomicity and Factorizations}

An element $a \in M \setminus M^\times$ is called an \emph{atom} if for all $b,c \in M$, the equality $a = bc$ implies that $M^\times \cap \{b,c\}$ is nonempty. We let $\mathcal{A}(M)$ denote the set consisting of all the atoms of $M$. An element $b \in M$ is called \emph{atomic} if either $b$ is a unit or $b$ factors into atoms (allowing repetitions). Following Cohn~\cite{pC68}, we say that the monoid~$M$ is \emph{atomic} if every element of~$M$ is atomic. In addition, $M$ is called \emph{strongly atomic} if $M$ is atomic and, for any $b,c \in M$, the subset $\{b,c\}$ of $M$ has an MCD.
\smallskip

An element $p \in M \setminus M^\times$ is called a \emph{prime} if for all $b,c \in M$ such that $p \mid_M bc$, either $p \mid_M b$ or $p \mid_M c$. As $M$ is assumed to be cancellative, every prime element of $M$ is an atom. A monoid is called an \emph{AP monoid} if every atom is prime. Following Coykendall, Dobbs, and Mullins~\cite{CDM99}, we say that $M$ is \emph{antimatter} if $\mathcal{A}(M)$ is the empty set. Observe that every antimatter monoid is trivially an AP monoid. Following Grams and Warner~\cite{GW75}, we say that~$M$ is an \emph{irreducible-divisor-finite monoid} if every nonunit of $M$ is divisible by an atom.
\smallskip

Assume for the rest of this section that $M$ is an atomic monoid. Let $\mathsf{Z}(M)$ denote the free commutative monoid on $\mathcal{A}(M_{\text{red}})$, and let $\pi \colon \mathsf{Z}(M) \to M_{\text{red}}$ denote the only monoid homomorphism that fixes every element of the set $M_{\text{red}}$. For each $b \in M$, we set
\[
  \mathsf{Z}(b) = \mathsf{Z}_M(b) = \pi^{-1} (bM^\times).
\]
An element $b \in M$ is called \emph{factorial} if $\mathsf{Z}_M(b)$ is a singleton, which means that $b$ has a unique factorization in $M$. Then $M$ is called a \emph{unique factorization monoid} (UFM) if every element of $M$ is factorial. It is well known that every UFM is a GCD monoid: indeed, a monoid is UFM if and only if it is an atomic GCD monoid.
\begin{figure}[ht]
  \begin{tikzcd} [cramped]
    \textbf{ GCD } \arrow[r, Rightarrow] \arrow[red, r, Leftarrow, "/"{anchor=center,sloped}, shift left=1.5ex] & \textbf{ pre-Schreier }  \arrow[r, Rightarrow] \arrow[red, r, Leftarrow, "/"{anchor=center,sloped}, shift left=1.5ex] &
    \textbf{ GL }  \arrow[r, Rightarrow] \arrow[red, r, Leftarrow, "/"{anchor=center,sloped}, shift left=1.5ex] &\textbf{AP}
  \end{tikzcd}
  \caption{Three properties that generalize the GCD property.}
  \label{fig:generalized-gcd diagram}
\end{figure}



\medskip
%%%%%%%%%%%%%%%%%%%%%%%%%%%%%
\subsection{Integral Domains}

A commutative ring with identity is called an \emph{integral domain} if its only zero-divisor is $0$. 
\smallskip

Let $R$ be an integral domain. We let $R^\times$ denote the group of units of $R$, while we let $R^*$ denote the multiplicative monoid consisting of all nonzero elements of~$R$. Observe that the nonzero principal ideals of $R$ can be obtained by joinning $\{0\}$ to the principal ideals of the monoid $R^*$, whence we see that $R$ satisfies the ACC on principal ideals if and only if $R^*$ satisfies the ACCP, in which case, we say that $R$ satisfies the \emph{ACCP}. Notice also that an integral domain $R$ is a UFD if and only if $R^*$ is a UFM.
\smallskip

Atomic and divisibility properties also tranfer naturally from $R^*$ to $R$ as eveyr nonzero element of $R$ is an atom/prime in the monoid $R^*$ if and only if it is so in the integral domain $R$. For any $r,s \in R^*$ with $r \mid_{R^*} s$ we write $r \mid_R s$. In general, we can adapt every monoidal property to the setting of integral domains through their multiplicative monoids. For instance, we say that an integral domain $R$ is an \emph{MCD domain} (resp., \emph{MCD-finite domain}) provided that $R^*$ is an MCD monoid (resp., MCD-finite monoid). We can do the same for the rest of the monoidal properties we have previously defined to define the notions of a q-GCD domain, a pre-Schreier domain, a GL domain, and an AP domain, an atomic domain, a strongly atomic, and IDF domain.
%\smallskip



\medskip
%%%%%%%%%%%%%%%%%%%%%%%%%%%%
\subsection{Monoid Algebras} 

For the rest of this section, let $M$ be a monoid, which is written additively. For an indeterminate $x$, the commutative ring consisting of all polynomial expressions in~$x$ with exponents in $M$ and coefficients in $R$ is called the \emph{monoid algebra} (or \emph{monoid domain}) of~$M$ over~$R$. Following the notation in~\cite{rG84}, we will denote the monoid algebra of~$M$ over $R$ by $R[x;M]$, although we simply write $R[M]$ if we see no danger of ambiguity. If $M$ is torsion-free (in addition to being cancellative and commutative) then $R[M]$ is an integral domain with
\[
  R[M]^\times = \{ux^m : u \in R^\times \text{ and } m \in M^\times \}
\]
(see~\cite[Theorems~8.1 and 11.1]{rG84}). We assume, for the rest of this section, that $M$ is a torsion-free monoid and, therefore, that $R[M]$ is an integral domain. 
\smallskip

We say that $M$ is a \emph{linearly ordered monoid} with respect to a given total order relation $\preceq$ if for all $b,c,d \in M$ the inequality $b \prec c$ implies that $b+d \prec c+d$. It follows from Levi's theorem~\cite{fL13} that a commutative semigroup $S$ with an identity element can be turned into a linearly ordered monoid if and only if $S$ is both cancellative and torsion-free. Thus, as we are assuming that $M$ is torsion-free, we can assume the existence of a total order $\preceq$ on~$M$ that is compatible with its operation.
\smallskip

Let $M$ be a linearly ordered monoid under a total order relation, which we denote by~$\preceq$. Then we can write any nonzero polynomial expression $f$ in the monoid algebra $R[M]$ as follows:
\[
  f := c_1 x^{m_1} + \dots + c_k x^{m_k}
\]
for some nonzero $c_1, \dots, c_k \in R$ and exponents $m_1, \dots, m_k \in M$ with $m_1 \succ \dots \succ m_k$. In this case, we call $\text{supp} \, f := \{m_1, \dots, m_k\}$ the \emph{support} of $f$, while we refer to $\deg f := m_1$ and $\text{ord} \, f := m_k$ as the \emph{degree} and \emph{order} of $f$, respectively. Notice that when the monoid $M$ is $\nn_0$ under the usual order, we recover the standard notions of degree, order, and support for standard polynomials in~$R[x]$. 
\smallskip

It is well known that a monoid algebra $R[M]$ is a GCD domain if and only if $R$ is a GCD domain and $M$ is a GCD monoid. It follows from~\cite[Proposition 8.3]{GP74} that when a monoid algebra $R[M]$ is an integral domain, $\dim R[M] \ge 1 + \dim R$. Moreover, when $R$ is a Noetherian domain and $M$ is a torsion-free monoid, it follows from~\cite[Corollary~2]{jO88} that $\dim R[M] = \dim R + \text{rank} \, M$, whence $R[M]$ has Krull dimension~$1$ if and only if~$R$ is a field and $M$ is a rank-$1$ torsion-free monoid (every rank-$1$ torsion-free monoid is isomorphic to an additive submonoid of $\qq$; see, for instance,~\cite[Section~24]{lF70}). The class of one-dimensional monoid algebras is central to this paper.
\smallskip

Background information on monoid algebras with emphasis on the ascent of algebraic properties from $(M,R)$ to $R[M]$, including the most significant progress on the same subject until 1984, can be found in Gilmer's book~\cite{rG84}.










\bigskip
%%%%%%%%%%%%%%%%%%%%%%%%%%
%%%%%%%%%%%%%%%%%%%%%%%%%%
\section{The MCD Property}
\label{sec:the MCD property}

This section is devoted to investigating the existence of MCDs, giving special attention to the class of rank-$1$ torsion-free monoids.

\medskip
%%%%%%%%%%%%%%%%%%%%%%%%%%%%
\subsection{Two Classes of \texorpdfstring{Rank-$1$}{Rank-1} MCD Monoids}

First, we identify two classes of MCD rank-$1$ monoids; one of these classes consists of atomic monoids while the other consists of antimatter monoids (i.e., monoids having no atoms).

\smallskip
%%%%%%%%%%%%%%%%%%%%%%%%%%
\subsubsection{A Class of Atomic MCD Monoids}

We start by identifying a class of atomic MCD rank-$1$ monoids. Our motivation comes from~\cite{GL22a} and~\cite{LWZ24}, where special subclasses were considered and proved to be atomic MCD monoids. Fix a pair $(P,k)$, where $P$ is an infinite set of primes and $k \in \nn$, and let $(p_n)_{n \ge 1}$ be the strictly increasing sequence with underlying set $P$. Throughout the first part of this section, we will consider the following rank-$1$ torsion-free monoid:
\begin{equation} \label{eq:main monoid 1}
  M_{P,k} := \Big\langle \tfrac{1}{p_j p_{j+2} \dots p_{j+2k}} : j \in \nn \Big\rangle \subseteq \qq_{\ge 0}.
\end{equation}
To be consistent with~\cite{LWZ24}, we call $M_{P,k}$ the \emph{prime reciprocal monoid} induced by the pair~$(P,k)$. One can readily verify that the set of atoms of $M_{P,k}$ is the defining generating set
\begin{equation} \label{eq:atoms of M_P}
  \mathcal{A}(M_{P,k}) = \Big\{ a_{j,k} :=  \frac1{p_j p_{j+2} \cdots p_{j+2k}} : j \in \nn \Big\}.
\end{equation}
On the other hand, $M_{P,k}$ does not satisfy the ACCP because $\big( \frac1{p_{2n}} + M_{P,k} \big)_{n \ge 1}$ is an ascending chain of principal ideals that does not stabilize. Indeed, for each $n \in \nn$, the identity
\[
  \frac1{p_{2n}} = \frac1{p_{2n+2}} + (p_{2n+2} - p_{2n})p_{2n+4} \cdots p_{2n+2k} \frac1{p_{2n} p_{2n+2} \dots p_{2n + 2k}}
\]
holds and, therefore, $\frac1{p_{2n+2}} \mid_{M_{P,k}} \frac1{p_{2n}}$. It turns out that the monoid $M_{P,k}$ satisfies the almost ACCP. Before proving this, it is convenient to show that every element of $M_{P,k}$ has a certain canonical sum decomposition (we assume that the sum over an empty index set is $0$).

\begin{lemma}[Canonical Sum Decomposition] \label{lem:CSD 1}
  For an infinite subset $P$ of $\pp$ and $k \in \nn$, let $M_{P,k}$ be as defined in~\eqref{eq:main monoid 1}. Then every element $r \in M_{P,k}^\bullet$ can be written uniquely as
  \begin{equation}\label{eq:gen-decomp}
    r = c_{k-1}(r) + \sum_{j=1}^{n(k)} c_j(r) a_{j,k}
  \end{equation}
  for an index $n(k) \in \nn_0$, an element $c_{k-1}(r) \in M_{P,k-1}$, and some coefficients $c_1(r), \dots, c_{n(k)}(r) \in \nn_0$ such that $c_j(r) \in \ldb 0, p_{j+2k} - 1 \rdb$ for every $j \in \ldb 1,n(k) \rdb$ while $c_{n(k)}(r) \ge 1$.
\end{lemma}

\begin{proof}
  Fix a nonzero element $r \in M_{P,k}$. As $r \in M_{P,k}$, every prime dividing the denominator of $r$ lies in~$P$. Let $p_\ell$ be the largest prime dividing the denominator of $r$ and then set $n(k) := \ell - 2k$. As $M_{P,k}$ is atomic with $\mathcal{A}(M_{P,k}) = \{a_{j,k} : j \in \nn\}$, we can certainly take an index $m(k) \in \nn_0$ such that
  \begin{equation}\label{eq:first-decomp}
    r = c'_{k-1} + \sum_{j=1}^{m(k)} c'_j a_{j,k}
  \end{equation}
  for some $c'_{k-1} \in M_{P,k-1}$ and $c'_1, \dots, c'_{m(k)} \in \nn_0$ such that $c'_{m(k)} > 0$ (for instance, we can take $c'_{k-1} = 0$). Among all such representations of $r$, choose one minimizing~$m(k)$. We first argue that $m(k) \le n(k)$.
  \smallskip

  \noindent \textsc{Claim.} $m(k) \le n(k)$.
  \smallskip

  \noindent \textsc{Proof of Claim.} Assume, toward a contradiction, that $m(k) > n(k)$. Since $m(k)+2k > n(k)+2k = \ell$, the maximality of~$\ell$ ensures that $p_{m(k) + 2k}$ does not divide the denominator of~$r$. Hence the last coefficient $c'_{m(k)}$ must be divisible by $p_{m(k)+2k}$ and, after writing $c'_{m(k)} = c p_{m(k)+2k}$ for some $c \in \nn$, we obtain that
  \[
    c'_{m(k)} a_{m(k),k} = c \frac{1}{p_{m(k)} p_{m(k)+2} \dots p_{m(k)+2(k-1)}} = c a_{m(k),k-1} \in M_{P,k-1}.
  \]
  Therefore, we can replace $c'_{m(k)}$ copies of the atom $a_{m(k),k}$ by $c$ copies of the element $a_{m(k),k-1}$, which is an atom in $M_{P,k-1}$, thereby producing another sum decomposition of~$r$ of the form in~\eqref{eq:first-decomp} but with a smaller largest index, which contradicts the minimality of $m(k)$. Hence $m(k) \le n(k)$, and the claim is established.
  \smallskip

  Next, we normalize the coefficients $c'_1, \dots, c'_{m(k)}$ using the following procedure, which we call \emph{normalization}.
  \smallskip

  \noindent \textsc{Normalization.} For each $j \in \ldb 1, m(k) \rdb$, write $c'_j = b_j p_{j+2k} + c_j$ for some $b_j, c_j \in \nn_0$ with $c_j \in \ldb 0, p_{j+2k} - 1 \rdb$ and note that
  \[
    c'_j a_{j,k} = (b_j p_{j+2k} + c_j) \frac1{p_j p_{j+2} \dots p_{j+2k}} = b_j \frac{1}{p_j p_{j+2} \dots p_{j+2(k-1)}} + c_j
    \frac{1}{p_j p_{j+2} \dots p_{j+2k}} = b_j a_{j,k-1} + c_j a_{j,k}.
  \]
  Therefore, we can write the left-hand side of~\eqref{eq:first-decomp} in the following way:
  \[
    r = c'_{k-1} + \sum_{j=1}^{m(k)} c'_j a_{j,k} = \bigg( c'_{k-1} + \sum_{j=1}^{m(k)} b_j a_{j,k-1} \bigg) + \sum_{j=1}^{m(k)} c_j a_{j,k}.
  \]
  The minimality of $m(k)$ ensures that $c_{m(k)} \neq 0$, for otherwise the term $c'_{m(k)}a_{m(k),k}$ could be absorbed into $M_{P,k-1}$ and we would obtain a representation of the form in~\eqref{eq:first-decomp} with a smaller largest index.
  Observe that $c_{k-1}(r) := c'_{k-1} + \sum_{j=1}^{m(k)} b_j a_{j,k-1} \in M_{P,k-1}$ and so, after setting $c_j(r) := c_j$ for every $j \in \ldb 1, m(k) \rdb$, we obtain the desired sum decomposition of $r$, and so the existence part of the proposition follows.
  \smallskip

  To show the uniqueness of the desired sum decomposition, let us consider two sum decompositions of $r$ satisfying the desired requirements: take indices $m,n \in \nn_0$, elements $c_{k-1}, c'_{k-1} \in M_{P,k-1}$, and coefficients $c_1, \dots, c_m, c'_1, \dots, c'_n \in \nn_0$ such that
  \begin{equation} \label{eq:CSD M_{P,k} uniqueness}
    c_{k-1} + \sum_{j=1}^{m} c_j a_{j,k} = r = c'_{k-1} + \sum_{j=1}^{n} c'_j a_{j,k},
  \end{equation}
  and $c_i \in \ldb 0, p_{i + 2k} - 1 \rdb$ and $c'_j \in \ldb 0, p_{j + 2k} - 1 \rdb$ for every $(i,j) \in \ldb 1, m \rdb \times \ldb 1,n \rdb$ and $c_m, c'_n \ge 1$. Assume without loss of generality that $m \le n$, and then apply the $p_{n+2k}$-adic valuation map to both sides of~\eqref{eq:CSD M_{P,k} uniqueness} to see that $m=n$. Now we can subtract both decompositions to obtain the equality $c'_{k-1} - c_{k-1} = \sum_{j=1}^m (c'_j - c_j) a_{j,k}$ and so, for each $j \in \ldb 1,m \rdb$, we can apply the $p_{j + 2k}$-adic valuation map to verify that $c_j = c'_j$. Therefore $c_{k-1} = c'_{k-1}$, and so both decompositions of $r$ are indeed the same, which concludes our proof of uniqueness.
\end{proof}

With the notation as in \eqref{eq:gen-decomp}, we say that $c_{k-1}(r)$ is the \emph{projection of} $r$ on $M_{P,k-1}$.

\begin{lemma} \label{lem:If projection on M{P,k-1} is zero then factorial}
  For an infinite subset $P$ of $\pp$ and $k \in \nn$, let $M_{P,k}$ be as defined in~\eqref{eq:main monoid 1}. An element of $M_{P,k}$ is factorial if its projection on $M_{P,k-1}$ is zero.
  %with projection on $M_{P,k-1}$ is fini
\end{lemma}

\begin{proof}
  Take $r \in M_{P,k}$ whose projection on the monoid $M_{P,k-1}$ is $0$. In light of Lemma~\ref{lem:CSD 1}, we can write $r = \sum_{j=1}^n c_j a_{j,k}$ for some $c_1, \dots, c_n \in \nn_0$ such that $c_j \in \ldb 0, p_{j + 2k} - 1 \rdb$ for every $j \in \ldb 1,n \rdb$. We can further assume that $c_n \ge 1$. Suppose that $\sum_{j=1}^m c'_j a_{j,k}$ is a factorization of $r$ in $M_{P,k}$. As the projection of $r$ on $M_{P,k-1}$ is $0$, the inequality $c'_j < p_{j + 2k}$ holds for every $j \in \ldb 1,m \rdb$ for otherwise, proceeding as in the existence part of the proof of Lemma~\ref{lem:CSD 1} we would obtain a positive projection of $r$ on $M_{P,k}$. Now, by completing with zero coefficients if necessary, we can assume that $m=n$. Finally, for each $j \in \ldb 1,n \rdb$, we can apply the $p_{j+2k}$-valuation map to both sides of the equality $0 = \sum_{j=1}^n (c'_j - c_j)a_{j,k}$ to deduce the equality $c_j = c'_j$ from the inequality $|c_j - c'_j| < p_{j+2k}$. Thus, we conclude that $r$ has only one factorization in $M_{P,k}$.
\end{proof}

Using the established canonical sum decomposition as a tool, we can now prove that $M_{P,k}$ is an MCD monoid, which generalizes~\cite[Theorem~5.3]{LWZ24}.

\begin{proposition} \label{prop:MPk-is-MCD}
  For an infinite subset $P$ of $\pp$ and $k \in \nn$, let $M_{P,k}$ be the monoid defined in~\eqref{eq:main monoid 1}. Then $M_{P,k}$ is an MCD monoid.
\end{proposition}

\begin{proof}
  We proceed by induction on the parameter $k \in \nn$. The base case $k=1$ was established in~\cite[Theorem~5.3]{LWZ24}. Thus, we assume that $k \ge 2$ and that $M_{P,k-1}$ is an MCD monoid. Fix a nonempty finite subset $S = \{s_1, \dots, s_t\} \subseteq M_{P,k}$, and let us prove that $S$ has an MCD in $M_{P,k}$. We may assume that every element of $S$ is nonzero because, otherwise, $0$ is an MCD of $S$. For each $s \in S$, write its canonical sum decomposition in $M_{P,k}$ as
  \begin{equation}\label{eq:canon-sr}
    s = c_{k-1}(s) + \sum_{j=1}^{n(s)} c_j(s) a_{j,k}
  \end{equation}
  for an index $n(s) \in \nn_0$, an element $c_{k-1}(s) \in M_{P,k-1}$, and some coefficients $c_1(s), \dots, c_{n(s)}(s) \in \nn_0$ such that $c_j(s) \in \ldb 0, p_{j+2k} - 1 \rdb$ for every $j \in \ldb 1,n(s) \rdb$ and $c_{n(s)}(s) \ge 1$.
  Now set
  \[
    m := \max\{n(s) : s \in S\},
  \]
  and set $c_j(s)=0$ whenever $s \in S$ and $j \in \ldb n(s)+1,m \rdb$. By the induction hypothesis, $M_{P,k-1}$ is an MCD monoid, so we can take
  \[
    d_{k-1} \in \mcd_{M_{P,k-1}} \{c_{k-1}(s) : s \in S\} \subseteq M_{P,k}.
  \]
  For each $j \in \ldb 1,m \rdb$, set $m_j := \min\{c_j(s) : s \in S\}$, and consider the element
  \[
    d_k := d_{k-1} + \sum_{j=1}^m m_j a_{j,k} \in M_{P,k}.
  \]
  For each $s \in S$, the element $d_{k-1}$ divides $c_{k-1}(s)$ in $M_{P,k-1}$, while $\sum_{j=1}^m m_j a_{j,k}$ divides $\sum_{j=1}^m c_j(s) a_{j,k}$ in $M_{P,k}$. Hence $d_k$ is a common divisor of~$S$ in $M_{P,k}$.

  We claim that the set of common divisors of $S-d_k$ in $M_{P,k}$ is finite. Let $e$ be one such common divisor. Since $d_k+e$ is a common divisor of $S$, the normalization procedure in Lemma~\ref{lem:CSD 1} shows that the projection of $d_k+e$ on $M_{P,k-1}$ is a common divisor of $\{c_{k-1}(s) : s \in S\}$. This projection is divisible by $d_{k-1}$, and the fact that $d_{k-1}$ is an MCD of $\{c_{k-1}(s) : s \in S\}$ forces it to be equal to $d_{k-1}$. Thus the projection of $e$ on $M_{P,k-1}$ is $0$, and no carry occurs when the canonical decompositions of $d_k$ and $e$ are added.

  If $e=0$ then there is nothing to prove. Therefore assume that $e$ is nonzero and write
  \[
    e = \sum_{j=1}^n b_j a_{j,k}
  \]
  in canonical form, where $b_j \in \ldb 0,p_{j+2k}-1\rdb$ for every $j \in \ldb 1,n\rdb$ and $b_n \ge 1$. For $s \in S$, set $q_j(s):=c_j(s)-m_j$ for $j \in \ldb 1,m\rdb$, and set $q_j(s):=0$ for $j>m$. If $b_j>0$ and $b_j>q_j(s)$ for every $s \in S$ then normalizing the equalities $s-d_k=e+(s-d_k-e)$ would produce a carry of one copy of $a_{j,k-1}$ in the projection of every element of $S-d_k$ on $M_{P,k-1}$. This would make $a_{j,k-1}$ a common divisor of $\{c_{k-1}(s)-d_{k-1} : s \in S\}$, contradicting that $d_{k-1}$ is an MCD of $\{c_{k-1}(s) : s \in S\}$. Therefore, for each $j$ with $b_j>0$, there exists $s \in S$ such that $b_j \le q_j(s)$. It follows that $j \le m$ and that $b_j \le \max\{q_j(s) : s \in S\}$ for every $j$. Hence there are only finitely many possibilities for $e$, proving the claim.

  Now choose a maximal element $d$ among the common divisors of $S-d_k$. If a nonzero element of $M_{P,k}$ divided every element of $S-(d_k+d)$ then adding it to $d$ would yield a common divisor of $S-d_k$ strictly larger than $d$, contradicting the maximality of $d$. Thus, the only common divisor of $S-(d_k+d)$ is $0$, and so $d_k+d$ is an MCD of~$S$ in $M_{P,k}$.
\end{proof}

\smallskip
%%%%%%%%%%%%%%%%%%%%%%%%%%%%%%%%%%%%%%%%%%%%%%%%%
\subsubsection{A Class of Antimatter MCD Monoids}

We now turn to exhibiting another class of non-atomic MCD rank-$1$ monoids. For each $p \in \pp$, the set $V_p := \nn_0\big[\frac1p\big]$ of $p$-adic nonnegative rationals is a rank-$1$ valuation monoid. Now for each subset $P$ of $\pp$, we set
\begin{equation} \label{eq:the GCD monoid M_P}
  M_P := \bigoplus_{p \in P} V_p.
\end{equation}
The class of monoids we explore in this section is $\{M_P : P \subseteq \pp \}$. As with the class exhibited earlier, the elements of each monoid $M_P$ have a natural sum decomposition, which is illustrated in the following proposition.

\begin{proposition}~\cite[Proposition 5.14]{CGGP25} \label{prop:canonical decomposition}
  Let $P$ be a subset consisting of primes, and consider the monoid $M_P := \bigoplus_{p \in P} V_p$. Each element $q \in M_P$ can be written uniquely as follows
  \begin{equation} \label{eq:canonical representation for internal sum of p-valuation}
    q = c_0 + \sum_{(p,n) \in P \times \nn} c_{p,n} \frac1{p^n},
  \end{equation}
  where $c_0 \in \nn_0$ and, for each $p \in P$, the sequence $(c_{p,n})_{n \ge 1}$ consists of nonnegative integer coefficients, almost all of which are zero, such that $0 \le c_{p,n} < p$ for every $n \in \nn$.
\end{proposition}

For each pair $(p,n) \in P \times \nn$, we let $c_0, c_{p,n} \colon M_P \to \nn_0$ be the functions defined as follows: $c_0(q) = c_0$ and $c_{p,n}(q) = c_{p,n}$, where $c_0$ and $c_{p,n}$ are the coefficients in the canonical sum decomposition of $q$ in~\eqref{eq:canonical representation for internal sum of p-valuation}. Now we introduce the following terminology.

\begin{definition}
  Let $P$ be a nonempty set consisting of primes, and let $M_P$ be the monoid in~\eqref{eq:the GCD monoid M_P}.
  \begin{itemize}
    \item For each $q \in M_P$, we call the right-hand side of~\eqref{eq:canonical representation for internal sum of p-valuation} the \emph{canonical sum decomposition} of $q$ in $M_P$.
      \smallskip

    \item For each pair $(p,n) \in P \times \nn_0$, we call the function $c_{p,n} \colon M_P \to \nn_0$ the $(p,n)$-\emph{canonical projection} of $M_P$.
  \end{itemize}
\end{definition}

It was also proved in~\cite[Proposition~5.6]{CGGP25} that the canonical $(p,n)$-projections just defined on the monoid $M_P$ satisfy the following basic properties.
\begin{enumerate}
  \item For each $p \in P$, the inequality $v_p(q) \ge 0$ holds if and only if $c_{p,n}(q) = 0$ for every $n \in \nn$.
    \smallskip

  \item For each $(p,n) \in P \times \nn$, the equality $c_{p,n}(q+r) = c_{p,n}(q) + c_{p,n}(r)$ holds for all $q,r \in M_P$.
    \smallskip

  \item $c_0(q+r) = c_0(q) + c_0(r)$ for all $q,r \in M_P$ such that $\gcd(\mathsf{d}(q), \mathsf{d}(r)) = 1$.
\end{enumerate}

Fix $q \in M_P$. It will be convenient in our proof that $M_P$ is an MCD monoid to compact the canonical sum decompositions of $q$ inside $M_P$ by setting, for each $p \in P$,
\[
  c_p(q) := \sum_{n \in \nn} c_{p,n}(q)\frac1{p^n} \in V_p,
\]
and then rewrite
the canonical sum decomposition of $q$ as follows:
\begin{equation}
    q = c_0(q) + \sum_{p \in P} c_p(q).
\end{equation}
As the following proposition indicates, canonical sum decompositions inside $M_P$ behave well with respect to divisibility.
\smallskip

\begin{proposition}~\cite[Proposition~5.17]{CGGP25} \label{prop:CSD and divisibility}
  Let $P$ be a nonempty set of primes, and let $M_P$ be the monoid defined in~\eqref{eq:the GCD monoid M_P}. For any $r,s \in M_P$, the following statements hold:
  \begin{enumerate}
    \item  $c_0(r)$ is the largest integer dividing $r$ in $M_P$.
      \smallskip

    \item If $r \mid_{M_P} s$ then $c_0(r) \le c_0(s)$.
      \smallskip

    \item If $r \mid_{M_P} s$ and $c_p(r) > c_p(s)$ then $c_0(r) < c_0(s)$ for every $p \in P$.
  \end{enumerate}
\end{proposition}

We are in a position to argue that $M_P$ is an MCD monoid for every nonempty set $P$ consisting of primes.

\begin{theorem} \label{thm:M_P is an MCD monoid}
  For any nonempty set $P$ consisting of primes, the monoid $M_P$ defined in~\eqref{eq:the GCD monoid M_P} is an MCD monoid.
\end{theorem}

\begin{proof}
  Fix a nonempty finite subset $S$ of $M_P$ and let us prove that $S$ has an MCD inside $M_P$. If $0 \in S$ then $0$ is the only common divisor of $S$ in $M_P$, and so $0$ is an MCD of $S$. Thus, from now on we can assume that $S \subseteq M_P^\bullet$.
  For each $s \in S$, write its compact canonical sum decomposition:
  \[
    s = c_0(s) + \sum_{p \in P} c_p(s),
  \]
  where $c_p(s) = \sum_{n \in \nn} c_{p,n}(s)/p^n$ for each $p \in P$. Now set
  \[
    C := \min\{c_0(s) : s \in S\} \quad \text{ and } \quad S_0 := \{s \in S : c_0(s) = C \}.
  \]
  Observe that $C$ is a common divisor of $S$, and part~(2) of Proposition~\ref{prop:CSD and divisibility} ensures that no common divisor of $S$ has integer part larger than $C$, that is, $c_0(d) \le C$ for every common divisor of $S$ in $M_P$. Now consider the set
  \[
    \mathcal{D} := \{d \in M_P : d \text{ is a common divisor of } S \text{ and } c_0(d)=C\}.
  \]
  This set is nonempty because $C \in \mathcal{D}$. Although $\mathcal{D}$ need not be finite, we now show that every element of $\mathcal{D}$ is dominated by an element in a finite subset of $\mathcal{D}$. Let
  \[
    F := \{p \in P : c_{p,n}(s) \ne 0 \text{ for some } s \in S \text{ and } n \in \nn\}.
  \]
  Observe that the set $F$ is finite. For each $p \in F$, set
  \[
    N_p := \max\{n \in \nn : c_{p,n}(s) \ne 0 \text{ for some } s \in S\}.
  \]
  For $p \in P \setminus F$, set $N_p:=0$ and consider the set
  \[
    \mathcal{D}_0 := \{d \in \mathcal{D} : c_{p,n}(d)=0 \text{ for all } p \in P \text{ and all } n>N_p\}.
  \]
  The set $\mathcal{D}_0$ is finite: if $d \in \mathcal{D}_0$ and $c_{p,n}(d) \ne 0$ then $p \in F$ and $1 \le n \le N_p$, while the canonical decomposition bounds the corresponding coefficient by
  \[
    1 \le c_{p,n}(d) \le p-1.
  \]
  Equivalently, every possible denominator of an element of $\mathcal{D}_0$ divides the integer $\prod_{p \in F} p^{N_p}$.

  Let us verify that every element of $\mathcal{D}$ is dominated by an element of $\mathcal{D}_0$. Fix $d \in \mathcal{D}$. If $c_{p,n}(d)>0$ for some $p \in P$ and some $n>N_p$, choose such an index $n$ maximal for the fixed prime $p$. For every $s \in S$, write $s=d+u_s$ with $u_s \in M_P$. Comparing the $p$-primary canonical expansions above the level $N_p$, and using the maximality of $n$, shows that the coefficient of $1/p^n$ in $u_s$ is nonzero; hence $u_s-1/p^n \in M_P$. Thus $d+1/p^n$ is still a common divisor of $S$ and still has integer part $C$, whence $d+1/p^n \in \mathcal{D}$ and $d \prec d+1/p^n$. Repeating this finite enlargement process eliminates the largest offending denominator for each prime after finitely many steps and creates no larger denominator. Therefore the process terminates at an element of $\mathcal{D}_0$ that dominates~$d$.

  Now consider the set $\mathcal{D}$ as a poset under the divisibility induced by $M_P$: for any $d_1, d_2 \in M_P$, we say that $d_1 \preceq d_2$ provided that $d_1 \mid_{M_P} d_2$. Since $\mathcal{D}_0$ is finite and every element of $\mathcal{D}$ is dominated by an element of $\mathcal{D}_0$, the poset $(\mathcal{D}, \preceq)$ has a maximal element. Note that, for every maximal element $m$ of $\mathcal{D}$ and every $s \in S_0$, it follows that $c_0(m) = C = c_0(s)$ and so $c_0(s-m)=0$.

  We proceed to prove that every maximal element in the poset $(\mathcal{D}, \preceq)$ is an MCD of $S$ in $M_P$. Suppose, toward a contradiction, that we can take an element $m \in \mathcal{D}$ that is maximal with respect to $\preceq$ but is not an MCD of $S$. Then there exists a nonzero element $e \in M_P$ such that $m+e$ is a common divisor of~$S$. Because~$e$ is a common divisor of $S_0-m$, for every $s \in S_0$, the relation $e \mid_{M_P} s-m$ holds, and so from the fact that $c_0(s-m) = 0$ we obtain that $c_0(e) = 0$ (in light of part~(2) of Proposition~\ref{prop:CSD and divisibility}). Since $m \mid_{M_P} m+e$, another use of part~(2) of Proposition~\ref{prop:CSD and divisibility} gives
  \[
    C = c_0(m) \le c_0(m+e).
  \]
  On the other hand, as $m+e$ is a common divisor of $S$, the maximality of $C$ as an integer part of a common divisor of $S$ guarantees that $c_0(m+e) \le C$. Hence $c_0(m+e) = C$. Therefore $m+e$ is a common divisor of $S$ such that $c_0(m+e) = C$, whence $m+e \in \mathcal{D}$. However, this contradicts the maximality of~$m$ in $\mathcal{D}$ because the fact that $e\neq 0$ implies that $m \prec m+e$ (because $M_P$ is a reduced monoid). Thus, the only common divisor of $S-m$ must be $0$, and so $m$ is an MCD of $S$ in $M_P$.
\end{proof}




\smallskip
%%%%%%%%%%%%%%%%%%%%%%%%%%%%%%%%%%%%%%%%%%%%%%%%%%%%%%%%%%%%%%%%%
\subsection{On the Ascent of the MCD Property to Monoid Algebras}

Our next goal is to argue that the MCD property ascends to monoid algebras over the class of pre-Schreier monoids. We use the following lemma as a tool.

\begin{lemma}\cite[Proposition~6.2]{GP74} \label{lem:exponent sum condition}
Let $R$ be an integral domain, and let $M$ be a cancellative torsion-free monoid. For nonzero $f(x), g(x) \in R[M]$,
\[
  |\emph{supp} \, g(x)|\,\emph{supp} \, f(x)+\emph{supp} \, g(x) = (|\emph{supp} \, g(x)|-1)\,\emph{supp} \, f(x)+\emph{supp} \, f(x)g(x),
\]
where $nA$ denotes the $n$-fold sumset of $A$.
\end{lemma}

The following proposition is well known, and it can be proved by mimicking the proof given in~\cite[Theorem~1.1]{mZ87} in the setting of integral domains.

\begin{proposition} \label{prop:Schreier interpolation}
  Let $M$ be a pre-Schreier monoid. Then for any $d_1, \dots, d_k, s_1, \dots, s_\ell \in M$ with $d_i \mid_M s_j$ for every $(i,j) \in \ldb 1,k \rdb \times \ldb 1, \ell \rdb$, there exists $d \in M$ such that $d_i \mid_M d$ for every $i \in \ldb 1,k \rdb$ and $d \mid_M s_j$ for every $j \in \ldb 1, \ell \rdb$.
\end{proposition}

The following lemma is the last tool we need to prove the main theorem of this section.

\begin{lemma} \label{lem:primal singletons}
  Let $M$ be a pre-Schreier monoid, and let $B$ and $C$ be finite nonempty subsets of~$M$. For any common divisor $a \in M$ of $B+C$, we can write $a = b+c$ for some elements $b,c \in M$ such that $b$ is a common divisor of $B$ and $c$ is a common divisor of~$C$.
\end{lemma}

\begin{proof}
  Set $m:=|B|$, and let $b_1, \dots, b_m$ be the elements of~$B$. Assume first that $|C|=1$, and write $C=\{c\}$. Since $M$ is a pre-Schreier monoid, for each $i \in \ldb 1,m \rdb$, the divisibility relation $a \mid_M b_i+c$ allows us to write $a=b'_i+c'_i$ for some $b'_i,c'_i \in M$ such that $b'_i \mid_M b_i$ and $c'_i \mid_M c$. Because $c'_i \mid_M a$ and $c'_i \mid_M c$ for every $i \in \ldb 1,m \rdb$, Proposition~\ref{prop:Schreier interpolation} ensures the existence of an element $c' \in M$ such that $c'_i \mid_M c'$ for every $i \in \ldb 1,m \rdb$, while $c' \mid_M a$ and $c' \mid_M c$. Write $a=b'+c'$ for some $b' \in M$. For each $i \in \ldb 1,m \rdb$, take $t_i \in M$ with $c'=c'_i+t_i$. Then
  \[
    b'_i+c'_i=a=b'+c'=b'+c'_i+t_i,
  \]
  and cancellativity gives $b'_i=b'+t_i$. Thus $b' \mid_M b'_i \mid_M b_i$ for every $i \in \ldb 1,m \rdb$. Hence $a=b'+c'$ is the desired decomposition in the case $|C|=1$.

  Now assume that $n:=|C|\ge 2$, and let $c_1,\dots,c_n$ be the elements of~$C$. Applying the singleton case to $B$ and $\{c_j\}$ for each $j \in \ldb 1,n \rdb$, we obtain elements $b'_j,c'_j \in M$ such that $a=b'_j+c'_j$, where $b'_j$ is a common divisor of $B$ and $c'_j \mid_M c_j$. Since $b'_j \mid_M a$ and $b'_j \mid_M b_i$ for every $(i,j) \in \ldb 1,m \rdb \times \ldb 1,n \rdb$, Proposition~\ref{prop:Schreier interpolation} guarantees the existence of an element $b' \in M$ such that $b'_j \mid_M b'$ for every $j \in \ldb 1,n \rdb$, while $b' \mid_M a$ and $b' \mid_M b_i$ for every $i \in \ldb 1,m \rdb$. Write $a=b'+c'$ for some $c' \in M$. For each $j \in \ldb 1,n \rdb$, take $t_j \in M$ with $b'=b'_j+t_j$. Then
  \[
    b'_j+c'_j=a=b'+c'=b'_j+t_j+c',
  \]
  and cancellativity gives $c'_j=t_j+c'$. Therefore $c' \mid_M c'_j \mid_M c_j$ for every $j \in \ldb 1,n \rdb$. Thus $a=b'+c'$, where $b'$ and $c'$ are common divisors of $B$ and $C$, respectively.
\end{proof}

We are in a position to prove the main theorem of this section.

\begin{theorem}
  Let $F$ be a field, and let $M$ be a cancellative and torsion-free pre-Schreier monoid. If $M$ is an MCD monoid then $F[M]$ is an MCD domain.
\end{theorem}

\begin{proof}
  Let $G$ be the Grothendieck group of $M$. As $M$ is cancellative, we can identify~$M$ as a submonoid of~$G$. As $M$ is torsion-free, $G$ is a torsion-free abelian group. Thus, Levi's theorem allows us to pick a total order $\preceq$ on $G$ such that $G$ is a linearly ordered group with respect to~$\preceq$. Set $V:=G_{\succeq 0}$. Then $V$ is a GCD monoid, and it follows from~\cite[Theorem~14.5]{rG84} that $F[V]$ is a GCD domain.

  Let $s_1(x),\dots,s_k(x)$ be pairwise distinct nonzero elements of $F[M]$. Since $M$ is an MCD monoid, the finite nonempty subset
  \[
    E:=\bigcup_{i=1}^k \text{supp}\,s_i(x)
  \]
  has an MCD $r \in M$. For each $i \in \ldb 1,k \rdb$, set $s'_i(x):=x^{-r}s_i(x)$, and set $S':=\{s'_1(x),\dots,s'_k(x)\}$. It is enough to prove that $S'$ has an MCD in $F[M]$: if $m(x)$ is an MCD of $S'$ then $x^rm(x)$ is an MCD of $\{s_1(x),\dots,s_k(x)\}$. By the choice of $r$, the only common divisors in $M$ of the set $\bigcup_{i=1}^k \text{supp}\,s'_i(x)$ are units.

  For each $i \in \ldb 1,k \rdb$, let $\alpha_i \in F^*$ be the coefficient of the term of $s'_i(x)$ of minimum degree, let $o_i:=\ord\,s'_i(x)$, and set
  \[
    t_i(x):=\alpha_i^{-1}x^{-o_i}s'_i(x) \in 1+F[G_{\succ 0}].
  \]
  Since $F[V]$ is a GCD domain, the finite set $T:=\{t_1(x),\dots,t_k(x)\}$ has a GCD in $F[V]$. Choose such a GCD $g(x)$ and scale it so that its constant coefficient is $1$. Then $\ord\,g(x)=0$, and, for every $i \in \ldb 1,k \rdb$, we can write
  \[
    t_i(x)=g(x)h_i(x)
  \]
  for some $h_i(x) \in F[V]$ with $\ord\,h_i(x)=0$.

  Write
  \[
    g(x)=1+\sum_{j=1}^n \lambda_jx^{a_j-b_j},
  \]
  where $n \in \nn_0$, $\lambda_1,\dots,\lambda_n \in F^*$, and $a_1,b_1,\dots,a_n,b_n \in M$, with $a_j-b_j$ running through the nonzero exponents in $\text{supp}\,g(x)$. Set $b:=0$ if $n=0$ and $b:=b_1+\dots+b_n$ otherwise. Then $f(x):=x^bg(x)$ belongs to $F[M]$, and the term $x^b$ belongs to $\text{supp}\,f(x)$. Let $d \in M$ be an MCD of $\text{supp}\,f(x)$ in $M$. Since $d \mid_M b$, we can write $q:=b-d \in M$. Now set
  \[
    m(x):=x^{-d}f(x)=x^qg(x) \in F[M].
  \]
  Since $d$ is an MCD of $\text{supp}\,f(x)$, the only common divisors in $M$ of $\text{supp}\,m(x)$ are units.

  We first prove that $m(x)$ divides every element of $S'$ in $F[M]$. Fix $i \in \ldb 1,k \rdb$ and set $\ell_i:=|\text{supp}\,h_i(x)|$. Applying Lemma~\ref{lem:exponent sum condition} in $F[V]$ to the equality $x^qt_i(x)=m(x)h_i(x)$ gives
  \[
    \ell_i\,\text{supp}\,m(x)+\text{supp}\,h_i(x)
    = (\ell_i-1)\,\text{supp}\,m(x)+\{q\}+\text{supp}\,t_i(x).
  \]
  We claim that $x^{o_i-q}h_i(x) \in F[M]$. To prove this, fix $\eta \in \text{supp}\,h_i(x)$ and set $\gamma:=o_i+\eta \in G$. Since $o_i+\text{supp}\,t_i(x)=\text{supp}\,s'_i(x) \subseteq M$, the displayed equality shows that, for all $\mu_1,\dots,\mu_{\ell_i} \in \text{supp}\,m(x)$, there exists $\delta \in M$ such that
  \[
    \mu_1+\dots+\mu_{\ell_i}+\gamma=q+\delta
  \]
  in $G$. Write $\gamma=u-v$ with $u,v \in M$. Then $q+v$ is a common divisor in $M$ of $\ell_i\,\text{supp}\,m(x)+\{u\}$. Repeatedly applying Lemma~\ref{lem:primal singletons} yields
  \[
    q+v=e_1+\dots+e_{\ell_i}+u'
  \]
  where each $e_j$ is a common divisor of $\text{supp}\,m(x)$ in $M$ and $u' \mid_M u$. Each $e_j$ is a unit, and so $q+v$ is associate to $u'$. Hence $q+v \mid_M u$, which means that $\gamma-q \in M$. Thus $o_i-q+\eta \in M$ for every $\eta \in \text{supp}\,h_i(x)$, proving the claim.

  The claim gives
  \[
    s'_i(x)=\alpha_i x^{o_i}t_i(x)=\alpha_i m(x)\big(x^{o_i-q}h_i(x)\big)
  \]
  with $x^{o_i-q}h_i(x) \in F[M]$. Hence $m(x)$ is a common divisor of $S'$ in $F[M]$.

  It remains to prove that $m(x)$ is maximal. For each $i$, set $Q_i(x):=x^{o_i-q}h_i(x) \in F[M]$, so that $s'_i(x)=\alpha_i m(x)Q_i(x)$. Let $w(x) \in F[M]$ be a common divisor of $Q_1(x),\dots,Q_k(x)$. Write $\theta:=\ord\,w(x)$ and $w_0(x):=x^{-\theta}w(x) \in F[V]$. Since $w(x)$ divides each $Q_i(x)$ in $F[M]$, the element $w_0(x)$ divides each $h_i(x)$ in $F[V]$: indeed, the corresponding quotient lies in $F[G]$, and its order is $0$ because both $w_0(x)$ and $h_i(x)$ have order $0$. Because $g(x)$ is a GCD of $T$ in $F[V]$, the only common divisors of $h_1(x),\dots,h_k(x)$ in $F[V]$ are units. Therefore $w_0(x) \in F^*$, and so $w(x)$ is a monomial in $F[M]$.

  If $w(x)=\lambda x^\theta$ with $\lambda \in F^*$ then $x^\theta$ divides every $Q_i(x)$ in $F[M]$, and therefore $x^\theta$ divides every $s'_i(x)$ in $F[M]$. Hence $\theta$ is a common divisor in $M$ of $\bigcup_{i=1}^k \text{supp}\,s'_i(x)$, so $\theta$ is a unit of $M$. Thus $w(x)$ is a unit of $F[M]$. Therefore the only common divisors of $S'/m(x)$ are units, and $m(x)$ is an MCD of $S'$ in $F[M]$.
\end{proof}



\smallskip
%%%%%%%%%%%%%%%%%%%%%%%%%%%%%%%%%%%%%%%%%%%%%%%%%%%%%%%%%%%%%%%%%%%%%%
\subsection{On the Ascent of Atomicity to Monoid Algebras over Fields}

In \cite{mR93}, Roitman disproved the ascent of atomicity to polynomial extensions, but he identified the MCD property as a sufficient condition for atomicity to ascend to polynomial extensions: for any MCD domain $R$, the polynomial ring $R[x]$ is atomic provided that $R$ is atomic \cite[Proposition~1.1]{mR93}. 

Although the ascent of atomicity to monoid algebras over fields has been considered in several recent papers (see \cite{GR25} and references therein), it is still unknown whether atomicity ascends to monoid algebras over a given field when restricted to the class of MCD monoids. 

We are in a position to establish the last main result of this section.

\begin{theorem}
  For each $p \in \pp$, there exists a rank-$1$ torsion-free MCD monoid~$M$ that is atomic but whose monoid algebra $\ff_p[M]$ is not atomic.
\end{theorem}

\begin{proof}
  By~\cite[Theorem~5.6]{GGP25}, there exist an integer $q \in \nn_{\ge 2}$ with $\gcd(p,q)=1$ and a nonconstant polynomial $P_q(x) \in \ff_p[x]$ of degree at most~$2$ such that $P_q\big(x^{q^n}\big)$ is irreducible in $\ff_p[x]$ for every $n \in \nn$. Set $\ell_0:=0$, and let $(\ell_n)_{n \ge 1}$ be a strictly increasing sequence of positive integers such that $q^{\ell_n-2\ell_{n-1}} > 2p^{n+1}$ for every $n \in \nn$. For each $n \in \nn$, set
  \[
    a_n := \frac{p^n q^{\ell_n} - 1}{2p^{2n}q^{\ell_n}} \quad \text{ and } \quad  b_n := \frac{p^n q^{\ell_n} + 1}{2p^{2n}q^{\ell_n}}.
  \]
  In~\cite[Proposition~5.2]{GGP25}, the authors proved that the rank-$1$ torsion-free monoid $M_{p,q}$ generated by the set $A_{p,q}:=\{a_n,b_n : n \in \nn\}$ is atomic with set of atoms $A_{p,q}$ and satisfies
  \[
    \nn_0\bigg[\frac1p\bigg] \subseteq M_{p,q} \subseteq \zz\bigg[\frac1p, \frac1q\bigg].
  \]
  Set $M:=M_{p,q}$. We first argue that $M$ is an MCD monoid. For this, we prove the following claim.
  \smallskip

  \noindent \textsc{Claim 1.} For each $m \in M$, there exists $N \in \nn_0$ such that, for all integers $n>N$ and $h \in \nn_0$, the following are equivalent:
  \[
    ha_n \mid_M m \text{ or } hb_n \mid_M m, \qquad ha_n+hb_n \mid_M m.
  \]
  \smallskip

  \noindent \textsc{Proof of Claim 1.} Fix $m \in M$. Let $N \in \nn$ be such that $\nu_q(m)>-\ell_{N-1}$ and $q^{\ell_{N-1}}-1>m$. Let $n>N$ be a positive integer.

  Suppose, for the sake of contradiction, that there exists $h \in \nn$ such that $ha_n \mid_M m$ but $ha_n+hb_n \nmid_M m$. Thus, there exist an integer $L$ and an atomic decomposition
  \[
    m = \sum_{i=1}^L \alpha_i a_i + \sum_{i=1}^L \beta_i b_i,
  \]
  where $\alpha_i,\beta_i \in \nn_0$ for every $i \in \ldb 1,L \rdb$ and $\alpha_n \ne \beta_n$. Let $j \in \nn_{\ge n}$ be the largest integer at most $L$ such that $\alpha_j \ne \beta_j$. Then, for each $i \in \ldb 1,j-1 \rdb$, we have $\nu_q(\alpha_i a_i+\beta_i b_i) \ge -\ell_i \ge -\ell_{j-1}$. Similarly, for each $i \in \ldb j+1,L \rdb$, the maximality of $j$ gives $\alpha_i=\beta_i$, and so
  \[
    \nu_q(\alpha_i a_i+\beta_i b_i)=\nu_q\bigg(\alpha_i \cdot \frac{p^iq^{\ell_i}-1}{2p^{2i}q^{\ell_i}}+\alpha_i \cdot \frac{p^iq^{\ell_i}+1}{2p^{2i}q^{\ell_i}}\bigg)=\nu_q\bigg(\frac{\alpha_i}{p^i}\bigg) \ge 0 \ge -\ell_{j-1}.
  \]
  Therefore as $\nu_q(m)>-\ell_{N-1}\ge -\ell_{j-1}$, we must have $\nu_q(\alpha_j a_j+\beta_j b_j)\ge -\ell_{j-1}$. This means that $\alpha_j-\beta_j$ is divisible in $\zz$ by $q^{\ell_j-\ell_{j-1}}$. Since $\alpha_j \ne \beta_j$, it follows that $\max\{\alpha_j,\beta_j\}\ge q^{\ell_j-\ell_{j-1}}$. Assume that $\alpha_j\ge q^{\ell_j-\ell_{j-1}}$, as the other case is analogous. Then
  \[
    \alpha_j a_j \ge q^{\ell_j-\ell_{j-1}} \cdot \frac{p^jq^{\ell_j}-1}{2p^{2j}q^{\ell_j}} = \frac{p^jq^{\ell_j}-1}{2p^{2j}q^{\ell_{j-1}}} \ge q^{\ell_{j-1}}-1,
  \]
  using our bounds on the sequence $(\ell_n)_{n \ge 1}$. This contradicts $q^{\ell_{j-1}}-1>m$. Similarly, the case $hb_n \mid_M m$ and $ha_n+hb_n \nmid_M m$ is impossible. Hence Claim~1 is established.
  \smallskip

  We still need to prove the following claim.
  \smallskip
  
  \noindent \textsc{Claim 2.} For each $m \in M$, there exists $d \in \nn_0\big[1/p\big]$ such that $d \mid_M m$ and, for any $d' \in \nn_0\big[1/p\big]$, the inequality $d'>d$ implies that $d' \nmid_M m$.
  \smallskip

  \noindent \textsc{Proof of Claim 2.} If $m \in \nn_0\big[1/p\big]$ then we can simply take $d:=m$. Therefore we assume that $m \notin \nn_0\big[1/p\big]$. Take $N \in \nn$ such that $\nu_q(m)>-\ell_{N-1}$ and $q^{\ell_{N-1}}-1>m$. In light of Claim~1, we can write $m$ as
  \[
    m = u + \sum_{i=1}^N \left(\alpha_i a_i + \beta_i b_i\right),
  \]
  where $\alpha_i,\beta_i \in \nn_0$ for every $i \in \ldb 1,N \rdb$ and $u \in \nn_0\big[1/p\big]$, by taking an atomic decomposition of~$m$ and observing that $a_n+b_n \in \nn_0\big[1/p\big]$ for every $n>N$. Note that $u$ is uniquely determined by the $(2N)$-tuple $(\alpha_1,\beta_1,\dots,\alpha_N,\beta_N)$, and also that
  \[
    \alpha_i,\beta_i \le m \bigg( \frac{p^iq^{\ell_i}-1}{2p^{2i}q^{\ell_i}} \bigg)^{-1} < 2mp^{i+1}.
  \]
  This implies that there are finitely many $(2N+1)$-tuples $(\alpha_1,\beta_1,\dots,\alpha_N,\beta_N,u)$ where $\alpha_i,\beta_i \in \nn_0$ for $i \in \ldb 1,N \rdb$ and $u \in \nn_0\big[1/p\big]$ with
  \[
    m = u+\sum_{i=1}^N \left(\alpha_i a_i + \beta_i b_i \right).
  \]
  Let $d$ be the maximum value of $u$ across all such tuples. Now take $d' \in \nn_0\big[1/p\big]$ with $d'>d$ and assume, for the sake of contradiction, that $d' \mid_M m$. Since $m \notin \nn_0\big[1/p\big]$, we see that $\nu_q(m)<0$, whence
  \[
    \nu_q(m-d')=\min\{\nu_q(m),\nu_q(d')\}=\nu_q(m),
  \]
  so $\nu_q(m-d')>-\ell_{N-1}$. Also, $q^{\ell_{N-1}}-1>m>m-d'$. Thus, by the same argument as above, we can write
  \[
    m-d'=u'+\sum_{i=1}^N \left(\alpha_i'a_i + \beta_i'b_i\right),
  \]
  where $\alpha_i',\beta_i' \in \nn_0$ for every $i \in \ldb 1,N \rdb$ and $u' \in \nn_0\big[1/p\big]$. Therefore
  \[
    m=(d'+u')+\sum_{i=1}^N \left(\alpha_i'a_i + \beta_i'b_i\right).
  \]
  However, as $d'+u'>d$ and $(\alpha_1',\beta_1',\dots,\alpha_N',\beta_N',d'+u')$ is a valid $(2N+1)$-tuple from above, this contradicts the maximality of~$d$. Hence $d' \nmid_M m$ for any $d'>d$, and Claim~2 is established.
  \smallskip

  We are in a position to prove that $M$ is an MCD monoid, fix a finite subset $S=\{s_1,\dots,s_n\}$ of~$M$. Take $N \in \nn$ such that $\nu_q(s)>-\ell_N$ and $q^{\ell_N}-1>s$ for every $s \in S$, and enlarge $N$, if necessary, so that Claim~2 applies to every $s \in S$ and every index greater than $N$. We construct sequences $c_1,c_2,\dots,c_N \in \nn_0$ and $d_1,d_2,\dots,d_N \in \nn_0$ satisfying the following conditions for each $i \in \ldb 1,N \rdb$:
  \begin{itemize}
    \item $c_i$ is the largest nonnegative integer such that $c_ia_i \mid_M s-\sum_{j=1}^{i-1}(c_ja_j+d_jb_j)$ for all $s \in S$, and
    \smallskip
    
    \item $d_i$ is the largest nonnegative integer such that $d_ib_i \mid_M s-\big(c_ia_i+\sum_{j=1}^{i-1}(c_ja_j+d_jb_j)\big)$ for all $s \in S$.
  \end{itemize}
  Now, for each $i \in \ldb 1,n \rdb$, define
  \[
    s_i' := s_i - \sum_{j=1}^N \left(c_ja_j+d_jb_j\right).
  \]
  Moreover, for each $i \in \ldb 1,n \rdb$, let $r_i$ denote the largest element of $\nn_0\big[1/p\big]$ dividing $s_i'$, whose existence follows from Claim~3. Choose $i_0 \in \ldb 1,n \rdb$ such that $r_{i_0}=\min\{r_1,r_2,\dots,r_n\}$, and consider the element
  \[
    d := r_{i_0}+\sum_{j=1}^N \left(c_ja_j+d_jb_j\right).
  \]
  We claim that $d$ is an MCD of~$S$. It suffices to show that no atom divides each element of $S-d$. It is convenient to split the rest of the proof into the following cases.
  \smallskip

  \textsc{Case 1:} $a_k$ is a common divisor of $S-d$ for some $k \le N$. Fix $s \in S$. As $a_k \mid_M s-d$, we have $a_k \mid_M s-d+m$ for any $m \in M$. Taking
  \[
    m = r_{i_0}+d_kb_k+\sum_{j=k+1}^N \left(c_ja_j+d_jb_j\right),
  \]
  we obtain
  \[
    a_k \mid_M s-\bigg( c_ka_k+\sum_{j=1}^{k-1}\left(c_ja_j+d_jb_j\right)\bigg).
  \]
  Hence $(c_k+1)a_k \mid_M s-\sum_{j=1}^{k-1}(c_ja_j+d_jb_j)$, contradicting the maximality of $c_k$.
  \smallskip

  \textsc{Case 2:} $b_k$ is a common divisor of $S-d$ for some $k \le N$. In this case, an argument similar to that used in Case~1 yields
  \[
    (d_k+1)b_k \mid_M s-\bigg(c_ka_k+\sum_{j=1}^{k-1}(c_ja_j+d_jb_j)\bigg)
  \]
  for every $s \in S$, contradicting the maximality of~$d_k$.
  \smallskip

  \textsc{Case 3:} $a_k$ is a common divisor of $S-d$ for some $k>N$. Fix $s \in S$ and note that
  \[
    \nu_q(d) \ge \min\left\{\nu_q(r_{i_0}),\nu_q(c_1a_1+d_1b_1),\dots,\nu_q(c_Na_N+d_Nb_N)\right\} \ge -\ell_N.
  \]
  Now,
  \[
    \nu_q(s-d) \ge \min\{\nu_q(s),\nu_q(d)\} \ge -\ell_N,
  \]
  and $q^{\ell_N}-1>s\ge s-d$. The argument in the proof of Claim~2 now shows that if $a_k \mid_M s-d$ then $a_k+b_k \mid_M s-d$ as well. However, $a_k+b_k=1/p^k \in \nn_0\big[1/p\big]$ and divides $s_{i_0}-d=s_{i_0}'-r_{i_0}$. Thus, $r_{i_0}+1/p^k \mid_M s_{i_0}'$, contradicting the maximality of $r_{i_0}$.
  \smallskip

  \textsc{Case 4:} $b_k$ is a common divisor of $S-d$ for some $k>N$. By similar reasoning, $r_{i_0}+1/p^k \mid_M s_{i_0}'$, again contradicting the maximality of $r_{i_0}$.
  \bigskip

  We have therefore proved that none of the atoms of $M$ is a common divisor of $S-d$ in $M$, so $d$ is a MCD of~$S$ in $M$. Hence $M$ is an MCD monoid.
  \smallskip

  It only remains to show that the monoid algebra $\ff_p[M]$ is not an atomic domain. For this, we first observe that $P_q(x^{1/p^k}) \in \ff_p[M]$ for all $k \in \nn_0$ because $\nn_0[1/p] \subseteq M$. Let $f$ be a factor of $P_q$ in $\ff_p[M]$, and write $P_q = fg$ for some $g \in \ff_p[M]$. Then we can pick $k \in \nn_0$ such that
  \[
    f\big(x^{(pq)^k}\big)g\big(x^{(pq)^k}\big) = P_q\big(x^{(pq)^k}\big) = P_q\big(x^{q^k}\big)^{p^k}
  \]
  in $\ff_p[x]$. Since $P_q\big(x^{q^k}\big)$ is irreducible in the UFD $\ff_p[x]$, we can write
  \[
    f\big(x^{(pq)^k}\big) = P_q\big(x^{q^k}\big)^t
  \]
  for some $t \in \nn_0$. After changing variables, one sees that $f = P_q\big(x^{1/p^k}\big)^t$. Thus, each nonunit factor of $P_q$ in $\ff_p[M]$ has the form $P_q\big(x^{1/p^k}\big)^t$ for some $k \in \nn_0$ and $t \in \nn$.

  Now assume, towards a contradiction, that $\ff_p[M]$ is atomic. Then $P_q = a_1 \cdots a_n$ for some irreducibles $a_1, \dots, a_n$ in $\ff_p[M]$. Since $a_1$ is a nonunit factor of $P_q$, we can pick $k \in \nn_0$ and $t \in \nn$ such that $a_1 = P_q\big(x^{1/p^k}\big)^t$. As $a_1$ is irreducible, the equality $t=1$ holds. However,
  \[
    a_1 = P_q\big(x^{1/p^k}\big) = P_q\big(x^{1/p^{k+1}}\big)^p,
  \]
  which contradicts the irreducibility of $a_1$ in $\ff_p[M]$. Hence $\ff_p[M]$ is not atomic.
\end{proof}










\bigskip
%%%%%%%%%%%%%%%%%%%%%%%%%%%%%%%%%
%%%%%%%%%%%%%%%%%%%%%%%%%%%%%%%%%
\section{The MCD-Finite Property}
\label{sec:the MCD-finite property}

This section is devoted to the study of the MCD-finite property.

\medskip
%%%%%%%%%%%%%%%%%%%%%%%%%%%%%
\subsection{A Class of Rank-1 MCD-finite Monoids}

We begin this section by proving that the monoids $M_{P,k}$ we introduced in~\eqref{eq:main monoid 1}, which we proved satisfy the MCD property, also satisfy the MCD-finite property. Recall that
\[
  M_{P,k} := \Big\langle \tfrac{1}{p_j p_{j+2} \dots p_{j+2k}} : j \in \nn \Big\rangle \subseteq \qq_{\ge 0},
\]
where $(p_n)_{n \ge 1}$ is a strictly increasing sequence of primes with underlying set $P$ and $k \in \nn$. We have also proved in Lemma~\ref{lem:CSD 1} that every nonzero $r \in M_{P,k}$ has a canonical sum decomposition as follows:
\[
  r = c_{k-1}(r) + \sum_{j=1}^{n(k)} c_j(r) a_{j,k}
\]
for an index $n(k) \in \nn_0$, an element $c_{k-1}(r) \in M_{P,k-1}$, and some coefficients $c_1(r), \dots, c_{n(k)}(r) \in \nn_0$ such that $c_j(r) \in \ldb 0, p_{j+2k} - 1 \rdb$ for every $j \in \ldb 1,n(k) \rdb$ while $c_{n(k)}(r) \ge 1$. Let us prove that every monoid $M_{P,k}$ is an MCD-finite monoid.

\begin{proposition} \label{prop:MPk-is-MCD-finite}
  For an infinite subset $P$ of $\pp$ and $k \in \nn$, let $M_{P,k}$ be the monoid defined in~\eqref{eq:main monoid 1}. Then $M_{P,k}$ is an MCD-finite monoid.
\end{proposition}

\begin{proof}
  We proceed by induction on the parameter $k \in \nn$. The base case $k=1$ follows from~\cite[Theorem~5.3]{LWZ24}, whose proof shows that every nonempty finite subset of $M_{P,1}$ has only finitely many MCDs. Assume that $k \ge 2$ and that $M_{P,k-1}$ is an MCD-finite monoid. Fix a nonempty finite subset $S := \{s_1, \dots, s_t\}$ of $M_{P,k}$. We will show that $S$ has only finitely many MCDs in $M_{P,k}$. As in the proof of Proposition~\ref{prop:MPk-is-MCD}, we may assume that every element of $S$ is nonzero. For each $s \in S$, write its canonical sum decomposition in $M_{P,k}$ as
  \[
    s = c_{k-1}(s) + \sum_{j=1}^{n(s)} c_j(s) a_{j,k}.
  \]
  Set $T := c_{k-1}(S) := \{c_{k-1}(s) : s \in S\}$. Then $T$ is a finite nonempty subset of $M_{P,k-1}$. We first prove the following claim.
  \smallskip

  \noindent \textsc{Claim.} If $d$ is an MCD of~$S$ in $M_{P,k}$ then $c_{k-1}(d)$ is an MCD of $T$ in $M_{P,k-1}$.
  \smallskip

  \noindent \textsc{Proof of Claim.} Let $d$ be an MCD of $S$ in $M_{P,k}$, and set $q := c_{k-1}(d)$. For each $s \in S$, apply the normalization procedure from the proof of the existence part of Lemma~\ref{lem:CSD 1} to the equality $d+(s-d)=s$. This shows that $q$ divides $c_{k-1}(s)$ in $M_{P,k-1}$. Hence $q$ is a common divisor of $T$.

  Now let $e \in M_{P,k-1}$ be a common divisor of $T-q$ in $M_{P,k-1}$. The same normalization argument shows that $e$ is a common divisor of $S-d$. Since $d$ is an MCD of $S$, the only common divisor of $S-d$ is $0$. Thus $e=0$, and so $q$ is an MCD of $T$ in $M_{P,k-1}$, as claimed.
  \smallskip

  By the induction hypothesis, the set $T$ has only finitely many MCDs in $M_{P,k-1}$. Thus it suffices to show that, for each fixed MCD $q$ of $T$ in $M_{P,k-1}$, the set
  \[
    \{ d \in M_{P,k} : d \in \mcd_{M_{P,k}}(S) \text{ and } c_{k-1}(d) = q \}
  \]
  is finite. Fix such an MCD $q$. If $d \in M_{P,k}$ is an MCD of $S$ satisfying $c_{k-1}(d)=q$ then $d=q+u$, where the projection of $u$ on $M_{P,k-1}$ is $0$. Also, $u$ is a common divisor of the finite set $S-q := \{s-q : s \in S\}$: indeed, $u=d-q \mid_{M_{P,k}} s-q$ for every $s \in S$. If a nonzero element of $M_{P,k}$ divided every element of $(S-q)-u$ then adding it to $d$ would yield a common divisor of $S$ strictly larger than $d$, contradicting that $d$ is an MCD. Thus, for fixed $q$, the possible elements $u$ are MCDs of $S-q$ whose projection on $M_{P,k-1}$ is $0$.

  It remains to observe that there are only finitely many such elements $u$. To see this, write
  \[
    u = \sum_{j=1}^n b_j a_{j,k}
  \]
  in canonical form for some $b_1, \dots, b_n$ in the discrete interval $\ldb 0,p_{j+2k}-1\rdb$ such that $b_n \ge 1$. Set $m := \max \{n(s):s \in S\}$, and put $c_j(s)=0$ whenever $j>n(s)$.

  Since $u$ divides every element of $S-q$, the normalization procedure in Lemma~\ref{lem:CSD 1}, applied to the equalities $s-q=u+(s-q-u)$, implies that, if $b_j>0$ and $b_j>c_j(s)$ for every $s\in S$ then the carry at $a_{j,k}$ contributes one copy of $a_{j,k-1}$ to every element of $T-q$. This would make $a_{j,k-1}$ a nonzero common divisor of $T-q$, contradicting that $q$ is an MCD of $T$. Hence, for every index $j$ with $b_j>0$, there exists an element $s\in S$ such that $b_j\le c_j(s)$. Therefore $j\le m$, and
  \[
    0 \le b_j \le \max \{c_j(s):s\in S\}
  \]
  for every $j$. Hence the index $n$ is bounded by $m$, and each coefficient $b_j$ can take only finitely many values. Therefore only finitely many elements $u$ can occur. It follows that, for the fixed projection $q$, there are only finitely many MCDs $d$ of $S$ satisfying $c_{k-1}(d)=q$.

  Combining this with the fact that $T$ has only finitely many MCDs in $M_{P,k-1}$, we conclude that $S$ has only finitely many MCDs in $M_{P,k}$. Therefore $M_{P,k}$ is an MCD-finite monoid.
\end{proof}


\medskip
%%%%%%%%%%%%%%%%%%%%%%%%%%%%%%%%%%%
\subsection{Ascent of the MCD-finite Property to Polynomial Rings}

The question of whether the IDF property ascends to polynomial extensions was posed by Anderson, Anderson, and Zafrullah~\cite{AAZ90} back in 1990. This remained open until Malcolmson and Okoh gave a negative answer to this question. It is well known that the IDF property is a generalization of the FF property. It turns out that the MCD-finite property does ascend to polynomial extensions.

\begin{theorem}
  If an integral domain $R$ is MCD-finite then $R[x]$ is also MCD-finite.
\end{theorem}

\begin{proof}
  Let $F$ be the field of fractions of $R$. Note that $F[x]$ is a Euclidean domain, so it is a UFD and every element of $F[x]$ has finitely many divisors up to associates. Since the units of $F[x]$ are precisely $F^\times$ and $a(x) \mid_{R[x]} b(x)$ implies $a(x) \mid_{F[x]} b(x)$, every element of $R[x]$ has finitely many divisors in $R[x]$ up to multiplication by an element of $F$.

  Assume for contradiction that there exists a set $S$ consisting of elements $f_k(x) \in R[x]$ for $k \in \ldb 1, n \rdb$ with infinitely many MCDs. Let $D$ denote the set of divisors of $f_1(x)$ in $F[x]$ such that no two elements of $D$ are associate, and let $\phi$ denote the map from MCDs $m(x)$ of $S$ to elements in $D$ such that $\phi(m(x)) \sim_{F[x]} m(x)$. By the pigeonhole principle, since there are finitely many elements in~$D$, there exist infinitely many MCDs of $S$ that are associate in $F[x]$.

  Let $a(x)$ be an MCD, and let $(r_m)_{m \ge 1}$ be a sequence of elements in $F$ such that the polynomials $r_m a(x)$ are also MCDs of $S$. For each $k \in \ldb 1, n \rdb$, define $b_k(x) \in R[x]$ such that $f_k(x) = a(x) b_k(x)$ and let $B \subseteq R$ be the union of the set of coefficients of $b_k(x)$. Since $a(x)$ is an MCD, there is no nonunit common divisor of $B$ in $R$ (otherwise $a(x) d$ would be a common divisor). Similarly, $B(r_m)^{-1} \subseteq R$ has no nonunit common divisor in $R$. Let $a_1$, $b_1$, and $f_1$ be the leading coefficients of $a(x)$, $b_1(x)$, and $f_1(x)$, respectively. Then we have $f_1 = a_1b_1$.

  We will now construct infinitely many MCDs of $f_1 B \subseteq R$, which contradicts that $R$ is MCD-finite. For all $m \in \nn$, we have $f_1 r_m = a_1b_1r_m = (a_1 r_m)b_1 \in R$ because $a_1 r_m$ is the leading coefficient of $r_m a(x) \in R[x]$. Also, we see that $f_1B (f_1 r_m)^{-1} = B(r_m)^{-1} \subseteq R$, and this subset has no nonunit common divisor in $R$; hence $f_1 r_m$ is an MCD of $f_1 B$. Since the $r_m$ are pairwise non-associate in $R$, the elements $f_1 r_m$ are not associate. Therefore there are infinitely many MCDs, as desired.
\end{proof}



\medskip
%%%%%%%%%%%%%%%%%%%%%%%%%%%%%%%
%%%%%%%%%%%%%%%%%%%%%%%%%%%%%%%
\subsection{The q-GCD Property}

We conclude this paper by considering the ascent of the q-GCD property, which is a specialization of the MCD-finite property. Recall that a cancellative commutative monoid is called a q-GCD monoid if every nonempty finite subset has at most one MCD, and an integral domain is called a q-GCD domain provided that its multiplicative monoid is a q-GCD monoid. Our primary purpose in this subsection is to argue that the q-GCD property does not ascend to monoid algebras over fields. First, we prove the following lemma.

\begin{lemma} \label{lem:CD Puiseux monoids are q-GCD}
  Let $M$ be an additive submonoid of $\qq_{\ge 0}$. If every nonempty finite subset of $M^\bullet$ has a common divisor in $M^\bullet$ then $M$ is a q-GCD monoid.
\end{lemma}

\begin{proof}
  Let $S = \{s_1, \dots, s_n\}$ be an arbitrary finite nonempty subset of $M$. We want to show that $|\mcd_M(S)| \le 1$. If $0 \in S$ then the only common divisor of $S$ is $0$, and so $0$ is the only MCD of $S$ in~$M$. Therefore we assume that $0 \notin S$. In addition, we can assume that $S$ has at least one MCD since otherwise $|\mcd_M(S)| \, \le 1$. Let $d$ be an MCD of $S$ in $M$, and let us argue the following claim.
  \smallskip

  \noindent \textsc{Claim.} $d \in S$.
  \smallskip

  \noindent \textsc{Proof of Claim.} Assume for the sake of contradiction that $d \notin S$. For each $i \in \ldb 1, n \rdb$, the inequality $s_i - d > 0$ holds because $d$ properly divides $s_i$ in $M$. Set
  \[
    S' = \big\{ s_i - d : i \in \ldb 1, n \rdb \big\}.
  \]
  By assumption, $S'$ has a nonzero common divisor $c \in M$. However, the fact that $d$ properly divides $c+d$ in $M$ contradicts the maximality of $d$ as a common divisor of $S$. Hence $d \in S$.
  \smallskip

  Because $d$ is a common divisor of $S$ that belongs to $S$, we obtain that $d = \min S$. Thus, the only MCD of a nonempty finite subset of $M$ must be its minimum, which completes the proof.
\end{proof}

It turns out that the q-GCD property does not ascend to monoid algebras over fields even in the class of rank-$1$ torsion-free monoids. We conclude with the following theorem, which proves this statement.

\begin{theorem}
  For every prime $p \in \pp$, there exists a rank-$1$ cancellative and torsion-free monoid $M$ such that $M$ is a q-GCD monoid but $\ff_p[M]$ is not a q-GCD domain.
\end{theorem}

\begin{proof}
  Fix $p \in \pp$. The first step of our proof is to construct an additive submonoid $M$ of the nonnegative cone of $\qq$ that is a q-GCD monoid. For this, pick an odd prime $q$ with $q \neq p$ and let~$V$ be the nonnegative cone of the additive abelian group $\zz[1/p]$, which is a valuation monoid. Here we denote the linearly ordered abelian group $\zz[1/p, 1/q]$ by $G$, and we let $W$ denote the nonnegative cone of $G$. Now consider the following additive monoid:
  \[
    M := V \cup G_{\ge 1},
  \]
  which is an intermediate monoid of the monoid extension $V \subset W$. Because $W$ is an overring of $M$, the monoid algebra $\ff_p[W]$ is an overring of $\ff_p[M]$. Note that $\ff_p[W]$ is a GCD domain because $W$ is a valuation monoid.
  \smallskip

  We can readily verify that every positive element of $M$ is divisible in $M$ by an element of the form~$1/p^n$ for some $n \in \nn_0$. Indeed, if a positive element $m \in M$ satisfies $m \le 1$ then $m \in V$ and so $m = a/p^n$ for some $(a,n) \in \nn \times \nn_0$, in which case $1/p^n \mid_M m$. On the other hand, if $m \in M_{>1}$ then we can take $n \in \nn$ large enough so that $1/p^n < m-1$ to obtain that $m - 1/p^n > 1$ and so the fact that $m - 1/p^n \in W$ ensures that $m - 1/p^n \in M_{>1} \subseteq M$, whence $1/p^n \mid_M m$. Therefore, for any finite nonempty subset $S$ of $M \setminus \{0\}$, we can pick $n \in \nn_0$ large enough so that $1/p^n$ is a common divisor of~$S$ in~$M$. Thus, it follows from Lemma~\ref{lem:CD Puiseux monoids are q-GCD} that $M$ is a q-GCD monoid.
  \smallskip

  Now we can focus on proving that the monoid algebra $\ff_p[M]$ is not a q-GCD domain. Let $k$ be a positive integer such that $p^{k-1} > q$, and consider the following polynomial expressions of $\ff_p[M]$:
  \[
    s_1(x) := x(1+x^{\frac{1}{pq}})^{p^k+1}(1+x^{\frac{1}{pq}} + x) \quad \text{ and } \quad s_2(x) := x^2(1+x^{\frac{1}{pq}})^2(1+x^{\frac{1}{pq}}+x).
  \]
  We will prove that the set $S:=\{s_1(x),s_2(x)\}$ has two non-associate MCDs in $\ff_p[M]$. Since the only units of $R$ are the nonzero elements of $\ff_p$, it suffices to show that
  \[
    g(x) := x(1+x^{\frac{1}{pq}})(1+x^{\frac{1}{pq}}+x) \quad \text{ and } \quad h(x) := x(1+x^{\frac{1}{pq}})^2
  \]
  are non-associate MCDs of $S$ in $\ff_p[M]$. Observe that $g(x), h(x) \in \ff_p[M]$ because $M_{\ge 1} = W_{\ge 1}$. It follows from the inequality $\deg g(x) > \deg h(x)$ that $g(x) \nmid_{\ff_p[M]} h(x)$. In addition, $1 + x^{1/pq} \nmid_{\ff_p[M]} 1 + x^{1/pq} + x$ since otherwise, we could take $A(x) \in \ff_p[M]$ such that $x = \big(1 + x^{1/pq}\big)\big(A(x)-1\big)$, whence $h(x) \nmid_{\ff_p[M]} g(x)$. Hence $g(x)$ and $h(x)$ are not associates in $\ff_p[M]$.
  \smallskip

  Next, we argue that $g(x)$ is an MCD of $S$ in $\ff_p[M]$. To this end, first observe that $g(x)$ is a common divisor of $S$ because
  \[
    \frac{s_1(x)}{g(x)} = \big( 1+x^{\frac1{pq}} \big)^{p^k} = 1 + x^{\frac{p^{k-1}}{q}} \in \ff_p[M] \quad \text{and} \quad \frac{s_2(x)}{g(x)} = x\big(1+x^{\frac1{pq}}\big) \in \ff_p[M],
  \]
  where the equality on the left follows from the Frobenius endomorphism and the first containment follows from the inequality $p^{k-1} > q$. To prove that $g(x)$ is maximal among the common divisors of $S$, observe that, inside the larger GCD domain~$\ff_p[W]$,
  \[
    \gcd\nolimits_{\ff_p[W]} \bigg( \frac{s_1(x)}{g(x)}, \frac{s_2(x)}{g(x)} \bigg) = \gcd\nolimits_{\ff_p[W]}\Big( 1 + x^{\frac{p^{k-1}}{q}}, \, x\big(1+x^{\frac{1}{pq}} \big)\Big) = 1 + x^{\frac{1}{pq}}.
  \]
  As $\ff_p[W]$ is a GCD domain, every common divisor of $S/g(x)$ in $\ff_p[M]$ must divide $1+x^{1/pq}$ in $\ff_p[W]$. Let $a(x) \in \ff_p[M]$ be such a common divisor. Since $a(x) \mid_{\ff_p[W]} 1+x^{1/pq}$, the largest exponent in $\operatorname{supp} a(x)$ is at most $1/pq$. Since $1/pq<1$, it follows that $\operatorname{supp} a(x) \subseteq V$. Take $\ell \in \nn$ such that $a\big(x^{p^\ell}\big) \in \ff_p[x]$. Because $a\big(x^{p^\ell q}\big)$ has integral exponents and nonzero constant term, the divisibility relation $a(x) \mid_{\ff_p[W]} 1+x^{1/pq}$ yields
  \[
    a\big(x^{p^\ell q}\big) \mid_{\ff_p[x]} 1+x^{p^{\ell-1}}=(1+x)^{p^{\ell-1}}.
  \]
  Because $\ff_p[x]$ is a UFD, there exist $u \in \ff_p^\times$ and $m \in \ldb 0,p^{\ell-1} \rdb$ such that $a\big(x^{p^\ell q}\big)=u(1+x)^m$. The left-hand side belongs to $\ff_p[x^q]$, while $(1+x)^m$ belongs to $\ff_p[x^q]$ only when $m=0$, as $q \neq p$. Hence $a\big(x^{p^\ell q}\big)$ is constant, and so $a(x)$ is constant. Thus $a(x) \in \ff_p[M]^\times$, and $g(x)$ is an MCD of $S$ in $\ff_p[M]$, as desired.

  Let us turn to verify that $h(x)$ is also an MCD of $S$ in $\ff_p[M]$. To argue that $h(x)$ is a common divisor of~$S$, we first notice that
  \[
    \frac{s_1(x)}{h(x)} = (1+x^{\frac1{pq}})^{p^k-1}(1+x^{\frac1{pq}}+x) = 1+x^{\frac{p^{k-1}}q} + x(1+x^{\frac1{pq}})^{p^k-1},
  \]
  and so the inequality $p^{k-1}/q > 1$ and the fact that the elements in the support of $x(1+x^{1/pq})^{p^k-1}$ are at least~$1$, ensures that $s_1(x)/h(x) \in \ff_p[M]$. In addition, note that all the exponents in the support of $s_2(x)/h(x) = x(1+x^{1/pq}+x)$ are at least~$1$, whence $s_2(x)/h(x) \in \ff_p[M]$. As a consequence, $h(x)$ is a common divisor of~$S$ in $\ff_p[M]$. In order to prove the maximality of $h(x)$. We first note that inside the GCD domain $\ff_p[W]$,
  \[
    \gcd\nolimits_{\ff_p[W]}\bigg( \frac{s_1(x)}{h(x)}, \frac{s_2(x)}{h(x)} \bigg) = 1 + x^{\frac1{pq}} + x,
  \]
  and so every common divisor of $S/h(x)$ in $\ff_p[M]$ divides $1+x^{1/pq}+x$ in $\ff_p[W]$. Let $d(x) \in \ff_p[M]$ be one of these common divisors. As $d(x)$ divides $1+x^{1/pq}+x$ in $\ff_p[W]$, the largest exponent in $\text{supp}\, d(x)$ is at most~$1$, whence $\text{supp}\, d(x) \subseteq \nn_0[1/p]$. This allows us to take $j \in \nn_0$ such that $d\big(x^{p^j}\big) \in \ff_p[x]$. From the divisibility relation $d(x) \mid_{\ff_p[W]} 1+x^{1/pq}+x$, we obtain
  \[
    d(x^{p^{j+1}q}) \mid_{\ff_p[x]} P(x) := x^{p^{j+1}q} + x^{p^j} + 1.
  \]
  We claim that $P(x) \in \ff_p[x]$ has no nonconstant divisors in $\ff_p[x^q]$. Assume, towards a contradiction, that $e(x^q) \in \ff_p[x^q]$ is a nonconstant divisor of $P(x)$ in $\ff_p[x]$. Since $e(x^q)$ is nonconstant, it has a root~$\alpha$ in some extension field of $\ff_p$. Then $P(\alpha)=0$, and $\alpha \neq 0$ because $P(0)=1$. Let $\zeta$ be a primitive $q$-th root of unity. Since $e((\alpha\zeta)^q)=e(\alpha^q)=0$, the element $\alpha\zeta$ is also a root of $P(x)$. Hence
  \begin{align*}
    \alpha^{p^j}(\zeta^{p^j}-1)
      &= \big((\alpha\zeta)^{p^{j+1}q}+ (\alpha\zeta)^{p^j}+1\big) - \big(\alpha^{p^{j+1}q}+\alpha^{p^j}+1\big) \\
      &= P(\alpha\zeta)-P(\alpha) = 0.
  \end{align*}
  Since $\alpha \neq 0$, it follows that $\zeta^{p^j}=1$, which is impossible because $\zeta$ has order $q$ and $q \nmid p^j$. Therefore $P(x)$ has no nonconstant divisors in $\ff_p[x^q]$. Since $d(x^{p^{j+1}q}) \in \ff_p[x^q]$ divides $P(x)$ in $\ff_p[x]$, the polynomial $d(x^{p^{j+1}q})$ is constant, and so $d(x)$ is constant.

  Hence $S$ has two non-associate MCDs in $\ff_p[M]$, whence the monoid algebra $\ff_p[M]$ is not a q-GCD domain.
\end{proof}






\bigskip
%%%%%%%%%%%%%%%%%%%%%%%%%%
%%%%%%%%%%%%%%%%%%%%%%%%%%
\section*{Acknowledgments}

During the preparation of this paper, the authors were part of PRIMES-USA at MIT, and they would like to thank the program for making possible the rewarding research experience that led to this paper. The second author kindly acknowledges support from the NSF under the award DMS-2213323.






\bigskip
%%%%%%%%%%%%%%%
%%%%%%%%%%%%%%%
\begin{thebibliography}{20}

\bibitem{ABLST23} K. Ajran, J. Bringas, B. Li, E. Singer, and M. Tirador, \emph{Factorization in additive monoids of evaluation polynomial semirings}, Comm. Algebra \textbf{51} (2023) 4347--4362.

\bibitem{dA00} D.~D. Anderson, \emph{GCD Domains, Gauss' Lemma, and Contents of Polynomials}. In: Non-Noetherian Commutative Ring Theory (Eds. S. T. Chapman and S. Glaz). Mathematics and Its Applications, vol. 520. Springer, Boston, MA.

\bibitem{AAZ90} D.~D. Anderson, D.~F. Anderson, and M. Zafrullah, \emph{Factorizations in integral domains}, J. Pure Appl. Algebra \textbf{69} (1990) 1--19.

\bibitem{BGLZ24} A. Bu, F. Gotti, B. Li, and A. Zhao, \emph{One-dimensional monoid algebras and ascending chains of principal ideals}. Submitted. Preprint on arXiv: \url{https://arxiv.org/abs/2409.00580}.

\bibitem{CGG20} S.~T. Chapman, F. Gotti, and M. Gotti, \emph{Factorization invariants of Puiseux monoids generated by geometric sequences}, Comm. Algebra \textbf{48} (2020) 380--396.

\bibitem{CGGP25} S.~T. Chapman, F. Gotti, M. Gotti, and H. Polo, \emph{On three families of dense Puiseux monoids}. Proceedings of Palermo Special Session in Commutative Algebra and Factorization Theory (to appear). Preprint on arXiv: \url{https://arxiv.org/abs/1701.00058}.

\bibitem{CG25} T. Cheng and F. Gotti, \emph{On divisibility aspects of commutative monoids}. Preprint.

\bibitem{pC68} P. M. Cohn, \emph{Bezout rings and their subrings}, Proc. Cambridge Philos. Soc. \textbf{64} (1968) 251--264.

%\bibitem{CG22} J. Correa-Morris and F. Gotti, \emph{On the additive structure of algebraic valuations of polynomial semirings}, J. Pure Appl. Algebra \textbf{226} (2022) 107104.

\bibitem{CDM99} J. Coykendall, D. E. Dobbs, and B. Mullins, \emph{On integral domains with no atoms}, Comm. Algebra \textbf{27} (1999) 5813--5831.

\bibitem{CG24} J. Coykendall and F. Gotti, \emph{Atomicity in integral domains}. In: Recent Progress in Rings and Factorization Theory (Eds. M. Brešar, A. Geroldinger, B. Olberding, and D. Smertnig), Springer Nature, Switzerland, 2024.

\bibitem{CG19} J. Coykendall and F. Gotti, \emph{On the atomicity of monoid algebras}, J. Algebra \textbf{539} (2019) 138--151.

\bibitem{DGHLS25} J. Dani, F. Gotti, L. Hong, B. Li, and S. Schlessinger, \emph{On finitary power monoids of linearly orderable monoids}. Submitted. Preprint on arXiv: \url{https://arxiv.org/abs/2501.03407}.

\bibitem{DG26} J. Du and F. Gotti, \emph{An unrestricted notion of the finite factorization property}. Submitted. Preprint on arXiv: \url{https://arxiv.org/abs/2511.00691}.

\bibitem{EK18} S. Eftekhari and M. R. Khorsandi, \emph{MCD-finite domains and ascent of IDF-property in polynomial extensions}, Comm. Algebra \textbf{46} (2018) 3865--3872.

\bibitem{lF70} L. Fuchs, \emph{Infinite Abelian Groups I}, Academic Press, 1970.

\bibitem{cG1801} C. F. Gauss, \emph{Disquisitiones Arithmeticae.} (translated by Arthur A. Clarke), Yale University Press, New Haven, 1966 (Corrected ed. Springer, New York, 1986), Art. 42.

\bibitem{GG25} A. Geroldinger and F. Gotti, \emph{On monoid algebras having every nonempty subset of $\mathbb{N}_{\ge 2}$ as a length set}, Mediterr. J. Math. \textbf{22} (2025) 1--19.

\bibitem{GH06} A. Geroldinger and F. Halter-Koch, \emph{Non-unique Factorizations: Algebraic, Combinatorial and Analytic Theory}, Pure and Applied Mathematics, vol. 278, Chapman \& Hall/CRC, Boca Raton, 2006.

\bibitem{rG84} R. Gilmer, \emph{Commutative Semigroup Rings}, Chicago Lectures in Mathematics, The University of Chicago Press, 1984.

\bibitem{GP74} R. Gilmer and T. Parker, \emph{Divisibility properties in semigroup rings}, Michigan Math. J. \textbf{21} (1974) 65--86.

\bibitem{GP25} V. Gonzalez and I. Panpaliya, \emph{On GL-domains and the ascent of the IDF property}. Submitted. Preprint on arXiv: \url{https://arxiv.org/abs/2504.10722}.

\bibitem{GGP25} V. Gonzalez, F. Gotti, and I. Panpaliya, \emph{On the ascent of almost and quasi-atomicity to monoid semidomains}. Submitted. Preprint on arXiv: \url{https://arxiv.org/abs/2501.04990}.

\bibitem{fG19} F. Gotti, \emph{Increasing positive monoids of ordered fields are FF-monoids}, J. Algebra \textbf{518} (2019) 40--56.

\bibitem{fG22} F. Gotti, \emph{On semigroup algebras with rational exponents}, Comm. Algebra \textbf{50} (2022) 3--18.

\bibitem{GL22a} F. Gotti and B. Li, \emph{Divisibility and a weak ascending chain condition on principal ideals}. Preprint on arXiv: \url{https://arxiv.org/abs/2212.06213}.

%\bibitem{GL22} F. Gotti and B. Li, \emph{Divisibility in rings of integer-valued polynomials}, New York J. Math \textbf{28} (2022) 117--139.

%\bibitem{GO20} F. Gotti and C. O'Neill, \emph{The elasticity of Puiseux monoids}, J. Commut. Algebra \textbf{12} (2020) 319--331.

\bibitem{GP23} F. Gotti and H. Polo, \emph{On the arithmetic of polynomial semidomains}, Forum Math. \textbf{35} (2023) 1179--1197.

%\bibitem{GP23a} F. Gotti and H. Polo, \emph{On the subatomicity of polynomial semidomains}. In: Algebra and Polynomials: Algebraic, Number Theoretic, and Topological Aspects of Ring Theory (Eds. J. L. Chabert, M. Fontana, S. Frisch, S. Glaz, and K. Johnson) pp. 197--212, Springer Nature, Switzerland, 2023.\\

\bibitem{GR25} F. Gotti and H. Rabinovitz, \emph{On the ascent of atomicity to monoid algebras}, J. Algebra \textbf{663} (2025) 857--881.

\bibitem{GZ23} F. Gotti and M. Zafrullah, \emph{Integral domains and the IDF property}, J. Algebra \textbf{614} (2023) 564--591.

\bibitem{GW75} A. Grams and H. Warner, \emph{Irreducible divisors in domains of finite character}, Duke Math. J. \textbf{42} (1975) 271--284.

\bibitem{dH1890} D. Hilbert, \emph{\"Uber die Theorie der algebraischen Formen}, Mathematische Annalen, 1890.

%\bibitem{JLZ23} N. Jiang, B. Li, and S. Zhu, \emph{On the primality and elasticity of algebraic valuations of cyclic free semirings}, Internat. J. Algebra and Comput. \textbf{33} (2023) 197--210.

%\bibitem{fK99} F. Kainrath, \emph{Factorization in Krull monoids with infinite class group}, Colloq. Math. \textbf{80} (1999) 23--30.

\bibitem{iK70} I. Kaplansky, \emph{Commutative rings}, Allyn and Bacon, Boston, 1970.

\bibitem{LP20} S. L'Innocente and F. Point, \emph{Bézout domains and lattice-valued modules}, J. Pure Appl. Algebra \textbf{224} (2020) 444--467.

\bibitem{fL13} F. W. Levi, \emph{Arithmetische Gesetze im Gebiete diskreter Gruppen}, Rend. Circ. Mat. Palermo \textbf{35} (1913) 225--236.

\bibitem{LWZ24} E. Liang, A. Wang, and L. Zhong, \emph{On maximal common divisors in Puiseux monoids}. Submitted. Preprint on arXiv: \url{https://arxiv.org/abs/2410.09251}.

\bibitem{MO06} P. Malcolmson and F. Okoh, \emph{A class of integral domains between factorial domains and IDF-domains}, Houston J. Math. \textbf{32} (2006) 399--421.

\bibitem{MO09} P. Malcolmson and F. Okoh, \emph{Polynomial extensions of idf-domains and of idpf-domains}, Proc. Amer. Math. Soc. \textbf{137} (2009) 431--437.

\bibitem{jO88} J. Okni\'nski, \emph{Commutative monoid rings with Krull dimension}. In: Semigroups Theory and Applications (Eds. H. J\"urgensen, G. Lallement, and H. J. Weinert), pp. 251--259. Lecture Notes in Mathematics, vol. 1320. Springer, Berlin, Heidelberg, 1988.

\bibitem{mR93} M. Roitman, \emph{Polynomial extensions of atomic domains}, J. Pure Appl. Algebra \textbf{87} (1993) 187--199.

\bibitem{TS67} T. Tamura and J. Shafer, \emph{On power semigroups}, Math. Jap. \textbf{12} (1967) 25--32.

\bibitem{mZ87} M. Zafrullah, \emph{On a property of pre-Schreier domains}, Comm. Algebra \textbf{15} (1987) 1895--1920.

\bibitem{mZ17} M. Zafrullah, \emph{Question HD 1704}. \url{https://lohar.com/mithelpdesk/hd1704.pdf}

\end{thebibliography}

\end{document}
